% !TeX root = ../../Book2.tex
% 纲要标题：## 第6章　张量积、双模与标量变换
% 纲要位置：计划纲要.md，第 1242 行。

\chapter{张量积、双模与标量变换}
\label{b2:ch:06}
\BookChapterNumber{b2:ch:06}

\begin{chapterabstract}
本章从平衡双变量映射构造张量积，并证明其泛性质。纯张量表达、双模作用、结合同构、Hom–张量伴随和扩张标量均给出明确的映射说明；商模与局部化的计算揭示张量积的右正合性及其可能不保持单射的原因。
\end{chapterabstract}

\begin{learninggoals}
\begin{itemize}
\item 从双可加性和平衡关系构造张量积，用泛性质证明映射良定义，区分纯张量与一般张量。
\item 检查双模的作用方向，写出结合、单位及反环同构的显式公式。
\item 证明 Hom–张量伴随及其自然性，辨认限制标量的左伴随与右伴随。
\item 证明张量积的右正合性，计算商模、局部化及基域扩张，并说明单射可能失效的原因。
\end{itemize}
\end{learninggoals}

设想一种由整数模 \(M,N\) 中的元素对产生的运算，暂记为 \(m\otimes n\)：它对两个变量分别可加，并允许整数在两因子之间移动，即 \((am)\otimes n=m\otimes(an)\)。如果所得交换群由这些符号生成，那么在 \(M=\mathbb Z/2\mathbb Z\)、\(N=\mathbb Z/3\mathbb Z\) 时，每个符号都是 \(\bar1\otimes\bar1\) 的整数倍。第一因子的关系迫使 \(2(\bar1\otimes\bar1)=0\)，第二因子的关系迫使 \(3(\bar1\otimes\bar1)=0\)；由 \(3-2=1\)，这个生成元只能是零。因而两个非零模所产生的群竟为零。本章将用生成元与关系构造这样的群，称为张量积，并证明它确实满足上述规则。

\ProseParagraphBreak

这类运算同时保留两因子的加法关系和标量关系。若想从 \((m,n)\) 定义一个目标交换群中的值，必须检验这些关系是否使结果保持不变；检验通过，双变量规则才能延伸为群同态。张量积的泛性质将准确说明这一对应，并帮助我们判断张量遇到商、直和与复合映射时发生什么。

\ProseParagraphBreak

本章使用第二章的模同态与商模、第三章的自由模与正合列，以及第五章的函子和自然性。环均含单位元，模作用均幺性；除明确假定交换的部分外，左右作用不能互换。张量积的泛性质把满足平衡条件的公式唯一地延伸为同态；函子语言进一步说明这些延伸如何随对象和映射变化。

\ProseParagraphBreak
交换环上的张量积、右正合性与标量变换可对照松村英之\cite[Appendix A, Tensor products, pp. 266--268; Appendix A, Change of coefficient ring, p. 268]{Matsumura1989CommutativeRingTheory}。伴随的自然性定义及扩张标量的例子见 Riehl\cite[Definition 4.1.1, pp. 116--117; Example 4.1.10(xi), p. 120]{Riehl2016CategoryTheory}。本章在一般环上分别核对左右作用，不把交换环中的对称公式直接移用。

% !TeX root = ../../Book2.tex
% 纲要标题：### 6.1 张量积的构造
% 纲要位置：计划纲要.md，第 980 行。

\section{张量积的构造}
\label{b2:sec:0601}

两个向量相乘通常没有预先给定的含义，但我们可以要求一种新的乘法对加法分配，并允许标量从一个因子移到另一个因子。张量积正是满足这些要求的通用对象。这里允许标量环非交换：一个因子是右模，另一个因子是左模，所得对象首先只是交换群。分清这一点，才能知道哪些式子是定义允许的运算。

% 纲要要点（供目录核对）：
%
% * 平衡双线性映射
% * 右模与左模的张量积
% * 纯张量表达的不唯一性
%

\subsection{平衡双线性映射}
\label{b2:subsec:0601:01}

\begin{definition}\label{b2:def:balanced-map}
设 \(R\) 为含单位元的环，\(M\) 为右 \(R\) 模，\(N\) 为左 \(R\) 模，\(A\) 为交换群。若映射 \(b:M\times N\to A\) 对任意 \(m,m'\in M\)、\(n,n'\in N\) 和 \(r\in R\) 满足
\[
\begin{aligned}
b(m+m',n)&=b(m,n)+b(m',n),\\
b(m,n+n')&=b(m,n)+b(m,n'),\\
b(mr,n)&=b(m,rn).
\end{aligned}
\]
就称 \(b\) 为一个 \(R\) 上的\term[pingheng-yingshe]{平衡映射}[balanced map]。前两式称为双可加性，第三式称为平衡条件。
\end{definition}

本小节标题沿用“平衡双线性映射”的常用说法，但在一般环上，其确切含义是上述“双可加且平衡”。目标 \(A\) 没有给定 \(R\) 模结构，因而不能写 \(b(mr,n)=r\cdot b(m,n)\)。若 \(R\) 交换且目标也是 \(R\) 模，通常所谓 \(R\) 双线性还要求两个变量分别 \(R\) 线性；它与这里的定义须在相应的模结构建立后比较。

\ProseParagraphBreak
例如，取 \(M=R_R\)、\(N={}_RR\)，环乘法 \(b(x,y)=xy\) 是到加法群 \((R,+)\) 的平衡映射，因为 \((xr)y=x(ry)\)。它说明平衡条件来自乘法结合律，不要求 \(xy=yx\)。又如，在 \(\mathbb Z\) 上，任何双可加映射都平衡：对正整数重复使用加法，先得到 \(b(am,n)=a\cdot b(m,n)=b(m,an)\)，再由负元推广到所有整数。

\subsection{右模与左模的张量积}
\label{b2:subsec:0601:02}

把每个有序对 \((m,n)\) 暂时看成一个独立符号 \([m,n]\)，构造以这些符号为基的自由交换群
\[
F=\mathbb Z^{(M\times N)}.
\]
这是\cref{b2:thm:free-module-universal}在 \(R=\mathbb Z\) 时的自由模。令 \(H\) 为以下三类元素生成的子群：
\[
\begin{aligned}
&[m+m',n]-[m,n]-[m',n],\\
&[m,n+n']-[m,n]-[m,n'],\\
&[mr,n]-[m,rn].
\end{aligned}
\]
这些关系只规定双可加性与平衡性，没有预设额外的模结构。

\begin{definition}\label{b2:def:tensor-product}
设 \(R\) 是环，\(M\) 是右 \(R\) 模，\(N\) 是左 \(R\) 模。令 \(F=\mathbb Z^{(M\times N)}\) 为以符号 \([m,n]\)（\(m\in M,n\in N\)）为基的自由交换群，并令 \(H\le F\) 由全部 \([m+m',n]-[m,n]-[m',n]\)、\([m,n+n']-[m,n]-[m,n']\) 及 \([mr,n]-[m,rn]\) 生成，其中 \(m,m'\in M\)、\(n,n'\in N\)、\(r\in R\)。考虑交换群 \(F/H\) 与映射
\[
\tau:M\times N\longrightarrow F/H,\qquad
(m,n)\longmapsto[m,n]+H
\]
称这对对象为右 \(R\) 模 \(M\) 与左 \(R\) 模 \(N\) 在 \(R\) 上的\term[zhangliangji]{张量积}[tensor product]，记 \(F/H=M\otimes_R N\)。将 \([m,n]+H\) 记为 \(m\otimes n\)，称为\term[chunzhangliang]{纯张量}[pure tensor]。
\end{definition}
\symidx{\(M\otimes_R N\)：右模与左模在 \(R\) 上的张量积}

由商掉的关系，\(\tau\) 平衡。双可加性还给出
\[
0\otimes n=m\otimes0=0,\qquad
(-m)\otimes n=-(m\otimes n)=m\otimes(-n).
\]
每个张量可以写成有限和 \(\sum_{i=1}^t m_i\otimes n_i\)：自由交换群中的整数系数可以移入第一个因子，负系数也由上式处理。\(t=0\) 时表示零元。

\begin{theorem}[张量积的泛性质]\label{b2:thm:tensor-universal}
设 \(R\) 是环，\(M\) 是右 \(R\) 模，\(N\) 是左 \(R\) 模，\(\tau:M\times N\to M\otimes_RN\) 是纯张量映射。对任意交换群 \(A\) 及平衡映射 \(b:M\times N\to A\)，存在唯一群同态
\[
\widetilde b:M\otimes_R N\longrightarrow A,
\qquad \widetilde b(m\otimes n)=b(m,n).
\]
反过来，每个这样的群同态与 \(\tau\) 复合都给出平衡映射。因此，从张量积出发的群同态与平衡映射一一对应。具有这一性质的交换群及其指定平衡映射，在保持指定映射的唯一同构意义下唯一。
\end{theorem}
\begin{proof}
先由自由交换群的泛性质，把 \([m,n]\mapsto b(m,n)\) 唯一延伸为同态 \(F\to A\)。双可加性与平衡条件恰好说明 \(H\) 的每个生成元都映到零，所以由\cref{b2:prop:quotient-module}的商模泛性质，该同态唯一通过 \(F/H\) 分解。唯一性也可直接看出：所有纯张量生成 \(F/H\)，它们的像已由公式给定。

反过来，同态保持加法，而 \(\tau\) 满足三类关系，所以复合映射平衡。若另有 \((T,\tau')\) 满足同样的泛性质，分别将 \(\tau'\) 与 \(\tau\) 分解，得到 \(M\otimes_RN\to T\) 与 \(T\to M\otimes_RN\)。两复合在指定平衡映射上的效果均为恒等，由各自泛性质的唯一性，两复合就是恒等映射。这同时给出同构及其唯一性。
\end{proof}

张量积的泛性质是下面的分解。图按底集合理解，\(\tau,b\) 是平衡映射，虚线是唯一的交换群同态。
\[
\begin{tikzcd}[column sep=3em,row sep=2.5em]
M\times N\arrow[r,"\tau"]\arrow[dr,"b"']&M\otimes_RN\arrow[d,dashed,"\exists!\,\widetilde b"]\\
&A.
\end{tikzcd}
\]
利用泛性质定义映射时，不必先判断某个张量的两种有限和表达是否相同；应当先证明有序对上的公式平衡。这个步骤正是表达不唯一时的良定义依据。

\begin{proposition}\label{b2:prop:tensor-maps}
设 \(R\) 是环，\(f:M\to M'\) 为右 \(R\) 模同态，\(g:N\to N'\) 为左 \(R\) 模同态，则存在唯一群同态
\[
f\otimes_R g:M\otimes_R N\longrightarrow M'\otimes_R N',
\qquad m\otimes n\longmapsto f(m)\otimes g(n).
\]
这些同态保持恒等映射，并满足
\[
(f'\otimes g')\circ(f\otimes g)=(f'\circ f)\otimes(g'\circ g).
\]
固定一个变量后，此构造还保持另一个变量中同态的加法。
\end{proposition}
\begin{proof}
双可加性来自 \(f,g\) 的可加性；平衡性由
\[
f(mr)\otimes g(n)=f(m)r\otimes g(n)
=f(m)\otimes r\cdot g(n)=f(m)\otimes g(rn)
\]
保证。于是\cref{b2:thm:tensor-universal}给出所需同态。恒等式与复合公式在每个纯张量上成立，故在整个群上成立。固定 \(g\) 为恒等映射时，\((f_1+f_2)(m)\otimes n=f_1(m)\otimes n+f_2(m)\otimes n\)，这证明第一变量中的加法相容性；固定第一变量时同样由第二变量的分配律得到。
\end{proof}

\subsection{纯张量表达的非唯一性}
\label{b2:subsec:0601:03}

在 \(\mathbb Z\otimes_{\mathbb Z}\mathbb Z\) 中，\(2\otimes3=1\otimes6=6\otimes1\)。因此，纯张量相等不意味着对应的因子相等。更强的现象是：在
\((\mathbb Z/2\mathbb Z)\otimes_{\mathbb Z}(\mathbb Z/3\mathbb Z)\) 中，任意纯张量 \(t\) 都满足 \(2t=3t=0\)，故 \(t=3t-2t=0\)。两边的模都非零，张量积却为零。

\ProseParagraphBreak
另一方面，一般张量并不一定是一个纯张量。为严格辨认这一点，先给出有基时的计算方法。此处只把所得对象看成交换群；下一节将补上域上的标量作用。

\begin{proposition}\label{b2:prop:tensor-free-coordinate}
设 \(R\) 是环，\(M=\bigoplus_{i\in I}e_iR\) 为有指定基 \((e_i)_{i\in I}\) 的右自由 \(R\) 模。对任何左 \(R\) 模 \(N\)，有群同构
\[
M\otimes_RN\longrightarrow N^{(I)},\qquad
\left(\sum_i e_i r_i\right)\otimes n\longmapsto(r_i n)_{i\in I}.
\]
其逆映射为 \((n_i)\mapsto\sum_i e_i\otimes n_i\)。若 \(R=k\) 是域，\(V,W\) 分别有基 \((e_i)\)、\((f_j)\)，则每个张量唯一写成有限和 \(\sum_{i,j}a_{ij}e_i\otimes f_j\)。
\end{proposition}
\begin{proof}
由于每个 \(m\in M\) 只有有限多个非零坐标，所给有序对上的映射确实取值于直和。逐坐标加法保持双可加性，\((r_i r)n=r_i(rn)\) 给出平衡性，故得到张量积上的同态。反向有限求和定义同态。正向后反向把 \(m\otimes n\) 送到 \(\sum_i e_i\otimes r_i n=m\otimes n\)；反向后正向恢复每个坐标，故互逆。再将每个 \(n_i\) 按 \(W\) 的基展开，便得到最后的存在性与唯一性。
\end{proof}

取 \(V=W=k^2\)，令 \(e_1,e_2\) 为标准基。张量
\[
t=e_1\otimes e_1+e_2\otimes e_2
\]
不是纯张量。若 \(t=(a_1e_1+a_2e_2)\otimes(b_1e_1+b_2e_2)\)，由唯一坐标得
\(a_1b_1=a_2b_2=1\)、\(a_1b_2=a_2b_1=0\)。前两式迫使四个因子均非零，与后两式矛盾。张量积的元素是满足关系的有限和；把它们全部想象成一个有序对，会遗漏这种最基本的例子。

\begin{proposition}[张量积与直和]\label{b2:prop:tensor-direct-sums}
设 \(R\) 是环，\((M_i)_{i\in I}\) 为右 \(R\) 模族，\(N\) 为左 \(R\) 模。存在自然的交换群同构
\[
\left(\bigoplus_{i\in I}M_i\right)\otimes_RN
\ \cong\ \bigoplus_{i\in I}(M_i\otimes_RN),
\qquad (m_i)\otimes n\longmapsto(m_i\otimes n)_i.
\]
固定右模而在左模变量取任意直和，也有同样的结论。
\end{proposition}
\begin{proof}
\((m_i)\) 有限支集，所以右端属于直和。有序对上的公式双可加且平衡，因而诱导同态。设 \(\iota_i:M_i\to\bigoplus_iM_i\) 为包含映射；各同态 \(\iota_i\otimes\operatorname{id}_N\) 经直和的泛性质合成反向同态。两复合分别在 \((m_i)\otimes n\) 和每个直和分量的纯张量上为恒等，故互逆。第二变量的公式为 \(m\otimes(n_i)\mapsto(m\otimes n_i)_i\)；相同的有限支集检查与包含映射构造给出其逆。各公式只使用给定映射与有限加法，所以与模同态相容。
\end{proof}

% !TeX root = ../../Book2.tex
% 纲要标题：### 6.2 双模与结合
% 纲要位置：计划纲要.md，第 1250 行。

\section{双模与结合}
\label{b2:sec:0602}

式子 \(mr\otimes n=m\otimes rn\) 使用了张量号两侧相邻的作用。如果 \(M\) 还有一个左作用，或 \(N\) 还有一个右作用，这些没有参与平衡关系的作用可能保留下来。双模使这一想法成为严格的构造，也说明张量积为什么可以连续进行。

% 纲要要点（供目录核对）：
%
% * 双模及张量积上的剩余作用
% * 张量积结合律与单位同构
% * 反环的必要性
%

\subsection{双模与张量积的剩余作用}
\label{b2:subsec:0602:01}

\begin{definition}\label{b2:def:bimodule}
设 \(S,R\) 为含单位元的环，\(P\) 是交换群，并给定幺性左 \(S\) 模与幺性右 \(R\) 模结构。如果两种标量作用对任意 \(s\in S\)、\(p\in P\)、\(r\in R\) 满足
\[
(sp)r=s(pr)\qquad(s\in S,\ p\in P,\ r\in R).
\]
就称 \(P\) 为 \((S,R)\) \term[shuangmo]{双模}[bimodule]，记为 \({}_SP_R\)。左、右各自满足模的全部公理，而上式要求两种作用彼此相容。
\end{definition}
\symidx{\({}_SP_R\)：左 \(S\)、右 \(R\) 双模}

环 \(R\) 以左右乘法成为 \({}_RR_R\)。矩形矩阵空间 \(M_{a\times b}(k)\) 是 \(\bigl(M_a(k),M_b(k)\bigr)\) 双模：左乘与右乘保留矩阵尺寸，矩阵乘法的结合律保证相容性。若 \(R\) 交换，每个左 \(R\) 模都通过 \(pr=rp\) 成为 \((R,R)\) 双模。

\begin{proposition}[剩余作用]\label{b2:prop:tensor-bimodule}
设 \(S,R,T\) 是环，\({}_SM_R\)、\({}_RN_T\) 为双模，则 \(M\otimes_RN\) 有唯一的 \((S,T)\) 双模结构，满足
\[
s(m\otimes n)=(sm)\otimes n,\qquad
(m\otimes n)t=m\otimes(nt).
\]
若只给定其中一种额外作用，则只得到相应一侧的模结构。
\end{proposition}
\begin{proof}
固定 \(s\in S\)，映射 \((m,n)\mapsto(sm)\otimes n\) 双可加，且
\[
\bigl(s(mr)\bigr)\otimes n=\bigl((sm)r\bigr)\otimes n=(sm)\otimes(rn).
\]
其中第一步使用双模相容性。因此由\cref{b2:thm:tensor-universal}得到群自同态 \(L_s\)。在纯张量上核对，可得 \(L_{s+s'}=L_s+L_{s'}\)、\(L_{ss'}=L_sL_{s'}\) 及 \(L_1=\operatorname{id}\)；纯张量生成全群，所以这些等式处处成立，给出左模结构。

固定 \(t\in T\)，用 \((rn)t=r(nt)\) 同样得到 \(D_t(m\otimes n)=m\otimes nt\)。右模公理在生成元上化为 \((nt)t'=n(tt')\) 等式。最后
\[
\bigl(s(m\otimes n)\bigr)t=(sm)\otimes nt
=s\bigl((m\otimes n)t\bigr),
\]
从而左右相容。由于作用在纯张量上的值已经给定，唯一性成立。
\end{proof}

两侧模结构决定张量后留下什么作用。各行只列由已给作用自然得到的结构：
\[
\begin{aligned}
M_R\otimes_R{}_RN&\quad\text{是交换群},\\
{}_SM_R\otimes_R{}_RN&\quad\text{是左}\;S\text{-模},\\
M_R\otimes_R{}_RN_T&\quad\text{是右}\;T\text{-模},\\
{}_SM_R\otimes_R{}_RN_T&\quad\text{得到}\;{}_S(M\otimes_RN)_T.
\end{aligned}
\]
相邻的两个 \(R\) 作用由平衡关系联系起来，外侧的左 \(S\)、右 \(T\) 作用则由上面的公式保留。若 \(R\) 非交换，交换两个因子时还须同时改用反环，见\cref{b2:eq:tensor-opposite-swap}。

\ProseParagraphBreak
若 \(R\) 交换，以上构造给出
\[
r(m\otimes n)=(rm)\otimes n=m\otimes(rn).
\]
因此 \(M\otimes_RN\) 是 \(R\) 模，\(\tau\) 是通常意义的 \(R\) 双线性映射。此时，对任意 \(R\) 模 \(L\)，一个 \(R\) 双线性映射 \(M\times N\to L\) 对应唯一的 \(R\) 线性映射 \(M\otimes_RN\to L\)：由群版本泛性质先得到同态，再在纯张量上核对 \(\widetilde b(rt)=r\cdot\widetilde b(t)\)。反方向也是直接复合。这就是交换环上熟悉的张量积泛性质。

\subsection{张量积结合律与单位同构}
\label{b2:subsec:0602:02}

\begin{theorem}[结合与单位同构]\label{b2:thm:tensor-associativity-unit}
设 \(A,R,S,T\) 是环，\({}_AM_R\)、\({}_RN_S\)、\({}_SP_T\) 为双模，则存在自然的 \((A,T)\) 双模同构
\[
\alpha:(M\otimes_RN)\otimes_SP\longrightarrow
M\otimes_R(N\otimes_SP),
\quad (m\otimes n)\otimes p\longmapsto m\otimes(n\otimes p).
\]
另有自然同构
\[
\rho:M\otimes_RR\longrightarrow M,\quad m\otimes r\longmapsto mr,
\qquad
\lambda:R\otimes_RN\longrightarrow N,\quad r\otimes n\longmapsto rn.
\]
它们保留所有已给定的外侧作用。没有给定 \(A,T\) 时，结合公式仍为交换群同构。
\end{theorem}
\begin{proof}
由于 \(m\otimes n\) 的表达不唯一，不能只把公式写在一组三个因子上，就认定它已定义在迭代张量积上。要证明公式不依赖纯张量及有限和的表达，可以依次使用两次泛性质。固定 \(p\in P\)，公式 \((m,n)\mapsto m\otimes(n\otimes p)\) 双可加，并对 \(R\) 平衡，因为
\[
mr\otimes(n\otimes p)=m\otimes r(n\otimes p)
=m\otimes(rn\otimes p).
\]
故得到群同态 \(h_p:M\otimes_RN\to M\otimes_R(N\otimes_SP)\)。于是 \((x,p)\mapsto h_p(x)\) 对第一变量可加；对第二变量的可加性可以在 \(x=m\otimes n\) 上检查，再由生成性推出。它对 \(S\) 平衡，因为在这些生成元上
\[
h_p\bigl((m\otimes n)s\bigr)
=m\otimes(ns\otimes p)=m\otimes(n\otimes sp)=h_{sp}(m\otimes n).
\]
再应用一次泛性质，得到 \(\alpha\)。

反向从 \((m,n,p)\mapsto(m\otimes n)\otimes p\) 出发。先固定 \(m\)，用 \((ns)\otimes p=n\otimes sp\) 的关系对 \(S\) 取张量，再用 \((mr)\otimes n=m\otimes rn\) 的关系对 \(R\) 取张量，得到 \(\beta\)。两复合在所有三重纯张量上为恒等，故互逆。公式还给出 \(\alpha(ax)=a\cdot\alpha(x)\)、\(\alpha(xt)=\alpha(x)t\)，所以保留外侧作用。

右乘 \((m,r)\mapsto mr\) 双可加且平衡，故定义 \(\rho\)。其逆为 \(m\mapsto m\otimes1\)，因为
\(mr\otimes1=m\otimes r\) 且 \(m1=m\)。左乘版本的逆是 \(n\mapsto1\otimes n\)，由 \(1\otimes rn=r\otimes n\) 与 \(1n=n\) 得到。若对 \(M,N,P\) 分别作相容同态，两条可能的复合路径在三重纯张量上的值都是各因子像的相同排列；单位公式也有同样的相容性。这证明所称的自然性。
\end{proof}

结合同构给出一个不依赖选基的指定映射，两边的底集合可以不同。例如，对 \({}_AM_R\)、\({}_RN_S\)、\({}_SP_T\) 及左 \(T\) 模 \(Q\)，从
\[
((M\otimes_RN)\otimes_SP)\otimes_TQ
\quad\text{到}\quad
M\otimes_R\bigl(N\otimes_S(P\otimes_TQ)\bigr),
\]
既可以先在最外层应用结合同构，把 \(P,Q\) 放在同一括号内，也可以先对 \(M,N,P\) 应用结合同构，再逐步改变剩余括号。两条路线都将
\[
((m\otimes n)\otimes p)\otimes q
\quad\longmapsto\quad
m\otimes\bigl(n\otimes(p\otimes q)\bigr).
\]
这些四重纯张量生成起点，所以两条路线给出同一个同态。对任意有限个依次可张量的双模，同样由多重纯张量的生成性得到这一相容性：结合映射只改变括号，不改变因子次序。以后省略括号，均指这些指定同构的识别，并不把两种构造当成集合意义下的相等。

\ProseParagraphBreak
交换环上还可以交换两个因子：\(m\otimes n\mapsto n\otimes m\)。平衡关系的像为
\(n\otimes rm=rn\otimes m\)，所以公式良定义，再交换一次就是恒等。一般非交换环上，\(N\otimes_RM\) 连所需的右、左作用都未必具备。例如，取域 \(k\)、\(n\ge2\) 及 \(R=M_n(k)\)，行向量空间 \(U=k^{1\times n}\) 按矩阵右乘成为右 \(R\) 模，列向量空间 \(V=k^{n\times1}\) 按矩阵左乘成为左 \(R\) 模，因此 \(U\otimes_RV\) 有定义。但这些数据没有指定 \(V\) 的右 \(R\) 作用，也没有指定 \(U\) 的左 \(R\) 作用，直接写 \(V\otimes_RU\) 就不满足构造所需的侧别条件。

\subsection{反环与作用侧别}
\label{b2:subsec:0602:03}

\subsecref{b2:subsec:0201:01}定义了反环 \(R^{\mathrm{op}}\)：底加法群不变，乘法次序反转。将右 \(R\) 模视为左 \(R^{\mathrm{op}}\) 模，可以整理作用侧别，却不能抹去次序信息。

\begin{proposition}\label{b2:prop:bimodule-opposite-action}
设 \(S,R\) 是环，\(P\) 是左 \(S\) 模。在 \(P\) 上指定与左 \(S\) 作用相容的幺性右 \(R\) 作用，等价于指定保持单位元的环同态
\[
R^{\mathrm{op}}\longrightarrow\operatorname{End}_S(P),
\qquad r\longmapsto D_r,
\]
其中 \(D_r(p)=pr\)，自同态环的乘法是复合。
\end{proposition}
\begin{proof}
双模相容性说明 \(D_r\) 是左 \(S\) 线性映射。右模结合律给出
\[
(D_r\circ D_t)(p)=(pt)r=p(tr)=D_{tr}(p).
\]
所以它反转 \(R\) 的乘法，恰好保持 \(R^{\mathrm{op}}\) 的乘法；分配律与幺性分别给出可加性和保持单位元。

反过来，由所给环同态定义 \(pr=D_r(p)\)。环同态的可加性给出标量加法的分配律，各 \(D_r\) 的可加性给出向量加法的分配律；又有 \(D_tD_r=D_{rt}\)，故 \((pr)t=p(rt)\)。保持单位元给出 \(p1=p\)，而 \(S\) 线性给出 \((sp)r=s(pr)\)。这便是需要的双模。
\end{proof}

因此，一份与左 \(S\) 作用相容的右 \(R\) 作用，正是一份 \(R^{\mathrm{op}}\) 在 \(P\) 上的 \(S\) 线性表示：每个反环元素作用为 \(\operatorname{End}_S(P)\) 中的一个自同态。反环使右作用的结合律与自同态的复合次序一致。

\ProseParagraphBreak
例如，左正则模 \({}_RR\) 的每个自同态都由 \(f(1)=a\) 决定，因为 \(f(r)=ra\)。将 \(a\) 对应于右乘映射，得到
\[
\operatorname{End}_R({}_RR)\cong R^{\mathrm{op}}.
\]
取 \(R=M_2(k)\)，对矩阵单位 \(E_{12},E_{21}\) 有
\[
(D_{E_{12}}\circ D_{E_{21}})(1)=E_{21}E_{12}=E_{22},
\qquad D_{E_{12}E_{21}}(1)=E_{11}.
\]
二者不同，直接表明把反环漏写会得到错误的乘法公式。

\ProseParagraphBreak
仍有一个始终合法的交换因子公式：把 \(N\) 改看为右 \(R^{\mathrm{op}}\) 模，把 \(M\) 改看为左 \(R^{\mathrm{op}}\) 模，就得到
\begin{equation}\label{b2:eq:tensor-opposite-swap}
M\otimes_RN\cong N\otimes_{R^{\mathrm{op}}}M,
\qquad m\otimes n\longmapsto n\otimes m.
\end{equation}
事实上，原关系 \(mr\otimes n=m\otimes rn\) 变为目标中的
\(n\otimes mr=rn\otimes m\)，正是反环上的平衡关系。反向交换给出逆。这是作用方向随交换一起改变的同构，与无条件在原环上交换因子是两回事。

% !TeX root = ../../Book2.tex
% 纲要标题：### 6.3 Hom–张量伴随
% 纲要位置：计划纲要.md，第 1256 行。

\section{Hom–张量伴随}
\label{b2:sec:0603}

一个双变量规则既可以借助张量积写成一个线性映射，也可以在固定一个变量后，把它读成以线性映射为值的映射。Hom–张量伴随说明这两种表达互相确定，并与各变量上的同态相容。非交换环上还须写清每一种线性条件所使用的标量作用。

% 纲要要点（供目录核对）：
%
% * 伴随同构及双模侧别；有限自由模的 Hom–张量同构及坐标逆式；无限自由模中的共同分母限制
% * 自然性与映射的显式公式
% * 限制标量、扩张标量及余诱导
%

\subsection{伴随同构及双模侧别}
\label{b2:subsec:0603:01}

设 \({}_SP_R\) 是双模，\(N\) 是左 \(S\) 模。\(\operatorname{Hom}_S(P,N)\) 表示左 \(S\) 线性映射构成的交换群。它有左 \(R\) 作用
\begin{equation}\label{b2:eq:hom-left-action}
(rf)(p)\coloneqq f(pr).
\end{equation}
令 \(D_r:P\to P\)、\(p\mapsto pr\)，则 \(rf=f\circ D_r\)：标量通过预复合输入端的右乘映射作用于 Hom。\(P\) 的右 \(R\) 作用因而变成 Hom 的左 \(R\) 作用，\(N\) 本身不需要有 \(R\) 作用。右端确实仍为 \(S\) 线性：
\[
(rf)(sp)=f\bigl((sp)r\bigr)=f\bigl(s(pr)\bigr)=s\cdot f(pr).
\]
又有 \(\bigl(r_1(r_2f)\bigr)(p)=f\bigl((pr_1)r_2\bigr)=\bigl((r_1r_2)f\bigr)(p)\)，其余分配律与单位条件由逐点加法给出。

\ProseParagraphBreak
现在设 \(M\) 为左 \(R\) 模。考虑规则 \(B:P\times M\to N\)，它对 \(P\) 是左 \(S\) 线性的，对 \(M\) 可加，并满足
\[
B(pr,m)=B(p,rm).
\]
张量积的泛性质把这个规则写成 \(h(p\otimes m)=B(p,m)\)。另一方面，固定 \(m\) 后，\(p\mapsto B(p,m)\) 属于 \(\operatorname{Hom}_S(P,N)\)，所以同一个规则也给出
\[
g:M\longrightarrow\operatorname{Hom}_S(P,N),\qquad
g(m)(p)=B(p,m).
\]
平衡条件正变为 \(g(rm)(p)=g(m)(pr)=(rg(m))(p)\)，即 \(g\) 的左 \(R\) 线性。下面的伴随同构就是在这两种表达之间转换。

\begin{theorem}[Hom–张量伴随]\label{b2:thm:hom-tensor-adjunction}
设 \(R,S\) 是环，\({}_SP_R\) 为双模，\(M\) 为左 \(R\) 模，\(N\) 为左 \(S\) 模。这些侧别给出
\[
P\otimes_RM\in S\text{-}\mathbf{Mod},\qquad
\operatorname{Hom}_S(P,N)\in R\text{-}\mathbf{Mod},
\]
其中后一个模的左 \(R\) 作用为 \((rf)(p)=f(pr)\)。此时有交换群同构
\begin{equation}\label{b2:eq:hom-tensor-adjunction}
\Phi:\operatorname{Hom}_S(P\otimes_RM,N)
\longrightarrow
\operatorname{Hom}_R\bigl(M,\operatorname{Hom}_S(P,N)\bigr).
\end{equation}
它及其逆 \(\Psi\) 由下列公式唯一确定：
\[
\Phi(h)(m)(p)=h(p\otimes m),\qquad
\Psi(g)(p\otimes m)=g(m)(p).
\]
此同构关于双模变量 \(P\) 和左 \(R\) 模变量 \(M\) 反变，关于左 \(S\) 模变量 \(N\) 协变。
\end{theorem}
\begin{proof}
先取左端的 \(h\)。固定 \(m\) 后，\(p\mapsto h(p\otimes m)\) 为左 \(S\) 线性，所以 \(\Phi(h)(m)\in\operatorname{Hom}_S(P,N)\)。关于 \(m\) 的可加性由张量分配律给出，而
\[
\Phi(h)(rm)(p)=h(p\otimes rm)=h(pr\otimes m)
=\bigl(r\cdot\Phi(h)(m)\bigr)(p)
\]
证明 \(\Phi(h)\) 是左 \(R\) 线性映射。

再取右端的 \(g\)。公式 \((p,m)\mapsto g(m)(p)\) 双可加，而且
\[
g(rm)(p)=\bigl(r\cdot g(m)\bigr)(p)=g(m)(pr),
\]
故平衡，诱导唯一群同态 \(\Psi(g)\)。由每个 \(g(m)\) 的左 \(S\) 线性，
\[
\Psi(g)(sp\otimes m)=g(m)(sp)=s\cdot g(m)(p),
\]
所以 \(\Psi(g)\) 还左 \(S\) 线性。

把两个公式代入，\(\Psi\Phi(h)\) 与 \(h\) 在每个纯张量上相同；\(\Phi\Psi(g)(m)\) 与 \(g(m)\) 在每个 \(p\) 上相同，因此两映射互逆。逐点加法说明它们是群同态。自然性将在下一小节按三个变量逐一核对。
\end{proof}

不需要交换性，也不要求模自由或有限生成。这一结论只把“对 \(P\) 与 \(M\) 分别满足相应线性条件的公式”改写成两种同态。若固定 \(P\)，记
\[
F(M)=P\otimes_RM,\qquad G(N)=\operatorname{Hom}_S(P,N),
\]
则\eqref{b2:eq:hom-tensor-adjunction}形如
\(\operatorname{Hom}_S\bigl(F(M),N\bigr)\cong\operatorname{Hom}_R\bigl(M,G(N)\bigr)\)。一般地，设 \(F:\mathcal C\to\mathcal D\)、\(G:\mathcal D\to\mathcal C\) 是函子。若每对对象 \(M\in\mathcal C\)、\(N\in\mathcal D\) 都有关于两变量自然的双射 \(\operatorname{Hom}_{\mathcal D}(F(M),N)\cong\operatorname{Hom}_{\mathcal C}(M,G(N))\)，就称 \(F,G\) 构成\term[bansui]{伴随}[adjunction]；\(F\) 称为 \(G\) 的\term[zuobansui]{左伴随}[left adjoint]，\(G\) 称为 \(F\) 的\term[youbansui]{右伴随}[right adjoint]，记作 \(F\dashv G\)。
\symidx{\(F\dashv G\)：函子 \(F\) 左伴随于 \(G\)}
自然性要求：改变 \(M\) 或 \(N\) 时，先复合同态再作上述转换，与先转换再复合得到同一映射。因此伴随的数据包含所有这些双射之间的相容性，逐对象任意选取 Hom 集的双射还不够。

\ProseParagraphBreak
\cref{b2:prop:tensor-direct-sums}中的直和公式也与这一对应相连。直和是模范畴的余积，从 \(\bigoplus_iM_i\) 出发的同态由各 \(M_i\) 上的同态任意指定；伴随将它们分别转成从 \(P\otimes_RM_i\) 出发的同态，故 \(P\otimes_R(\bigoplus_iM_i)\) 仍具有这些张量积的余积性质。

\ProseParagraphBreak
对反环应用同一定理，再交换张量因子的次序，就得到右模记法；双变量规则与固定变量的方法没有改变。同一结论写为
\[
\operatorname{Hom}_{S^{\mathrm{op}}}(M\otimes_RP,N)
\cong\operatorname{Hom}_{R^{\mathrm{op}}}
\bigl(M,\operatorname{Hom}_{S^{\mathrm{op}}}(P,N)\bigr),
\]
这里 \(M_R,{}_RP_S,N_S\)，内层 Hom 的右 \(R\) 作用是 \((fr)(p)=f(rp)\)。下标 \(S^{\mathrm{op}}\) 表示右 \(S\) 线性，避免把两种 Hom 混为一谈。

\ProseParagraphBreak
伴随把两种 Hom 联系起来，并不直接把 Hom 写成张量积。在有限自由的情形，坐标泛函还给出下面的同构。

\begin{proposition}\label{b2:prop:finite-free-hom-tensor}
设 \(R\) 为交换环，\(P\) 为有限自由 \(R\) 模，\(M\) 为任意 \(R\) 模，记 \(P^*=\operatorname{Hom}_R(P,R)\)。则求值诱导自然的 \(R\) 模同构
\[
\Theta:P^*\otimes_RM\longrightarrow\operatorname{Hom}_R(P,M),
\qquad \lambda\otimes m\longmapsto\bigl(p\mapsto\lambda(p)m\bigr).
\]
\end{proposition}
\begin{proof}
交换性保证 \(p\mapsto\lambda(p)m\) 是 \(R\) 线性映射，而这一赋值对 \((\lambda,m)\) 双线性，因而诱导 \(\Theta\)。
选 \(P\) 的有限基 \(e_1,\ldots,e_n\) 及其坐标泛函 \(e_i^*\)，反向定义
\[
\Psi(g)=\sum_{i=1}^n e_i^*\otimes g(e_i).
\]
对 \(p\in P\)，基展开给出
\[
\Theta\Psi(g)(p)=\sum_i e_i^*(p)g(e_i)=g(p).
\]
另一方面，\(\lambda=\sum_i\lambda(e_i)e_i^*\)，所以
\[
\Psi\Theta(\lambda\otimes m)
=\sum_i e_i^*\otimes\lambda(e_i)m
=\left(\sum_i\lambda(e_i)e_i^*\right)\otimes m
=\lambda\otimes m.
\]
纯张量生成定义域，故两映射互逆。\(\Theta\) 与 \(M\) 上的后复合、\(P\) 上的前复合相容，因而是自然同构；它的逆虽用基写出，却由 \(\Theta\) 唯一确定。
\end{proof}
这里的有限基使逆公式确为有限张量和。有限自由只是保证这个求值映射可逆的一种充分条件；\cref{b2:prop:finite-projective-hom-tensor}将用有限对偶基把结论推广到有限生成投射模。对于无限自由模，张量中的有限和则可能无法给出所有同态，下面的例子准确说明了缺少哪些映射。

\begin{proposition}\label{b2:prop:infinite-free-hom-tensor}
令 \(P=\bigoplus_{n\geq1}\mathbb Ze_n\) 为自由交换群。自然的 \(\mathbb Q\) 线性映射
\[
\Theta:\operatorname{Hom}_{\mathbb Z}(P,\mathbb Z)
\otimes_{\mathbb Z}\mathbb Q
\longrightarrow\operatorname{Hom}_{\mathbb Z}(P,\mathbb Q),
\qquad
\Theta(f\otimes q)(p)=qf(p)
\]
的像恰由满足下述条件的同态 \(h\) 组成：存在一个正整数 \(d\)，使对每个 \(n\geq1\) 都有 \(dh(e_n)\in\mathbb Z\)。特别地，\(h(e_n)=2^{-n}\) 定义的同态不在此像中，所以 \(\Theta\) 不满射。
\end{proposition}
\begin{proof}
赋值 \((f,q)\mapsto[p\mapsto qf(p)]\) 双可加且对整数标量平衡，故确实诱导所写的映射。任一张量可写成有限和 \(\sum_{i=1}^s f_i\otimes q_i\)。选取这有限个有理数 \(q_i\) 的一个共同正整数分母 \(d\)，则 \(dq_i\in\mathbb Z\)，因而其像 \(h\) 满足
\[
dh(e_n)=\sum_{i=1}^s(dq_i)f_i(e_n)\in\mathbb Z
\]
对所有 \(n\) 同时成立。零张量可以取 \(d=1\)。

反过来，若 \(h\) 满足这个共同分母条件，指定 \(f(e_n)=dh(e_n)\) 就给出唯一同态 \(f:P\to\mathbb Z\)。这是因为 \(P\) 中每个元素仅有有限多个非零坐标；整数值可按基展开逐项延拓。此时 \(\Theta(f\otimes d^{-1})=h\)，证明了像的描述。

自由模的同一延拓性质允许任意指定有理数值，故 \(h(e_n)=2^{-n}\) 确实定义同态。对于任一正整数 \(d\)，选取 \(n\) 使 \(2^n>d\)，便有 \(0<dh(e_n)<1\)，不能是整数。因此这个 \(h\) 没有共同整数分母。
\end{proof}

\subsection{自然性与映射的显式公式}
\label{b2:subsec:0603:02}

\secref{b2:sec:0502}中的自然性要求映射变动时同构仍然相容。下面核对伴随同构对三个变量的这种相容性。伴随公式没有使用基；有限自由情形的 \(\Theta\) 也由求值直接定义，其坐标逆式不依赖所选基则是因为一个同构的逆唯一。

\ProseParagraphBreak
在 \(M\) 变量中，同态 \(a:M'\to M\) 通过前复合产生反向映射；在 \(N\) 变量中，同态 \(b:N\to N'\) 通过后复合产生同向映射。若 \(a\) 左 \(R\) 线性、\(b\) 左 \(S\) 线性，将这两个变化合在一起，自然性公式是
\begin{equation}\label{b2:eq:tensor-adjunction-naturality}
\Phi\bigl(b\circ h\circ(\operatorname{id}_P\otimes a)\bigr)
=\operatorname{Hom}_S(P,b)\circ\Phi(h)\circ a.
\end{equation}
左边在 \(m'\)、\(p\) 处的值为 \(b\Bigl(h\bigl(p\otimes a(m')\bigr)\Bigr)\)；右边按复合定义给出同一个值。因此该等式成立，既证明了 \(M\) 变量中的反变相容性，也证明了 \(N\) 变量中的协变相容性。

\ProseParagraphBreak
记 \(\Phi_{M,N}\) 为相应对象上的伴随同构，则同一个自然性等式给出下面的交换方块。这里 \(F(a)=\operatorname{id}_P\otimes a\)、\(G(b)=\operatorname{Hom}_S(P,b)\)，两条竖箭头分别在 Hom 群上作所标出的复合。
\[
\begin{tikzcd}[column sep=3.5em,row sep=3em]
\operatorname{Hom}_S(F(M),N)\arrow[r,"\Phi_{M,N}"]\arrow[d,"{b\circ(-)\circ F(a)}"']&\operatorname{Hom}_R(M,G(N))\arrow[d,"{G(b)\circ(-)\circ a}"]\\
\operatorname{Hom}_S(F(M'),N')\arrow[r,"\Phi_{M',N'}"']&\operatorname{Hom}_R(M',G(N')).
\end{tikzcd}
\]

\(P\) 变量也通过前复合反变。若 \(c:P'\to P\) 是 \((S,R)\) 双模同态，则
\[
\Phi\bigl(h\circ(c\otimes\operatorname{id}_M)\bigr)
=\operatorname{Hom}_S(c,N)\circ\Phi(h),
\]
因为在 \(m,p'\) 处两边都是 \(h\bigl(c(p')\otimes m\bigr)\)。这里 \(\operatorname{Hom}_S(c,N)\) 为前复合，它左 \(R\) 线性恰由 \(c(p'r)=c(p')r\) 保证。这完成了\cref{b2:thm:hom-tensor-adjunction}中全部自然性断言的证明。

\ProseParagraphBreak
伴随双射还将两侧的恒等态射变成一对结构映射。在 \(\operatorname{Hom}_S(F(M),F(M))\cong\operatorname{Hom}_R(M,GF(M))\) 中，恒等态射 \(\operatorname{id}_{F(M)}\) 在右侧对应
\[
\eta_M:M\longrightarrow\operatorname{Hom}_S(P,P\otimes_RM),
\qquad \eta_M(m)(p)=p\otimes m.
\]
同样，在 \(\operatorname{Hom}_R(G(N),G(N))\) 中取 \(\operatorname{id}_{G(N)}\)，经 \(\Psi\) 转到另一侧，得到求值映射
\[
\varepsilon_N:P\otimes_R\operatorname{Hom}_S(P,N)\longrightarrow N,
\qquad p\otimes f\longmapsto f(p).
\]
它们分别称为这一伴随的\term[bansui-danwei]{单位}[unit]与\term[bansui-yudanwei]{余单位}[counit]。前面的伴随同构已保证其良定义和线性。将这两个态射转回原来的一侧，互逆性应恢复相应的恒等态射；自然性把这一往返转换写成下面两个复合。这就是\term[sanjiao-hengdengshi]{三角恒等式}[triangle identities]：
\symidx{\(\eta_M,\varepsilon_N\)：Hom–张量伴随的单位与余单位}
\[
\varepsilon_{F(M)}\circ(\operatorname{id}_P\otimes\eta_M)
=\operatorname{id}_{F(M)},\qquad
\operatorname{Hom}_S(P,\varepsilon_N)\circ\eta_{G(N)}
=\operatorname{id}_{G(N)}.
\]
省略函子复合的括号，这两个恒等式也就是下面的交换三角形；各箭头均为已经构造的单位、余单位或其函子像。
\[
\begin{tikzcd}[column sep=4em,row sep=2.5em]
F(M)\arrow[r,"F(\eta_M)"]\arrow[dr,"\operatorname{id}"']&FGF(M)\arrow[d,"\varepsilon_{F(M)}"]\\
&F(M)
\end{tikzcd}
\qquad
\begin{tikzcd}[column sep=4em,row sep=2.5em]
G(N)\arrow[r,"\eta_{G(N)}"]\arrow[dr,"\operatorname{id}"']&GFG(N)\arrow[d,"G(\varepsilon_N)"]\\
&G(N).
\end{tikzcd}
\]
第一式将 \(p\otimes m\) 送到 \(\eta_M(m)(p)=p\otimes m\)；第二式将 \(f\) 送到映射 \(p\mapsto\varepsilon_N(p\otimes f)=f(p)\)。这两个三角恒等式成立时，\(\eta_M\) 与 \(\varepsilon_N\) 仍可能不是同构；要得到范畴等价，还须证明单位与余单位在每个对象上都可逆。

\subsection{限制标量、扩张标量及余诱导}
\label{b2:subsec:0603:03}

设 \(\varphi:R\to S\) 为保持单位元的环同态，不假定单射。限制标量只让 \(R\) 中的元素通过 \(\varphi\) 作用；反方向则有两种构造：张量积在保留原有 \(R\) 关系的同时加入 \(S\) 标量，Hom 则把从 \(S\) 出发的全部 \(R\) 线性映射组成一个 \(S\) 模。前一种构造将给出限制标量的左伴随，后一种给出右伴随。

\ProseParagraphBreak
对左 \(S\) 模 \(N\)，令 \(r\cdot n=\varphi(r)n\)，就得到左 \(R\) 模；这个函子称为\term[xianzhi-biaoliang]{限制标量}[restriction of scalars]，记为 \(\operatorname{Res}_{\varphi}\)。它保留底交换群与映射，只改变允许使用的标量。

\ProseParagraphBreak
令 \(S\) 通过左乘及 \(t\cdot r=t\varphi(r)\)（\(t\in S,r\in R\)）成为 \((S,R)\) 双模。对左 \(R\) 模 \(M\)，左 \(S\) 模
\[
S\otimes_RM
\]
称为从 \(R\) 到 \(S\) 的\term[kuozhang-biaoliang]{扩张标量}[extension of scalars]。映射 \(m\mapsto1\otimes m\) 左 \(R\) 线性，因为 \(1\otimes rm=\varphi(r)\otimes m\)，但它不保证单射。

\ProseParagraphBreak
在右伴随中，同一个 \(S\) 改由 \(r\cdot t=\varphi(r)t\) 接受左 \(R\) 作用，同时保留右 \(S\) 乘法。左右两种结构及其产生的 Hom 作用列于\cref{b2:tab:scalar-change-sides}；表中的 \(M\) 是左 \(R\) 模。
\begin{table}[htbp]
\centering
\caption[标量变换的作用侧别]{标量变换中同一个 \(S\) 的两种双模结构与 Hom 作用}
\label{b2:tab:scalar-change-sides}
\begin{tabularx}{\linewidth}{@{}p{0.30\linewidth}X p{0.21\linewidth}@{}}
\toprule
对象 & 标量作用 & 所用构造 \\
\midrule
\({}_SS_R\) & 左 \(S\)：\(s\cdot t=st\)；右 \(R\)：\(t\cdot r=t\varphi(r)\) & \(S\otimes_RM\) \\
\({}_RS_S\) & 左 \(R\)：\(r\cdot t=\varphi(r)t\)；右 \(S\)：\(t\cdot s=ts\) & \(\operatorname{Hom}_R(S,M)\) \\
\(\operatorname{Hom}_R({}_RS,M)\) & 左 \(S\)：\((sf)(t)=f(ts)\)；不使用右作用 & 余诱导 \\
\bottomrule
\end{tabularx}
\end{table}

\begin{proposition}[限制标量的左右伴随]\label{b2:prop:scalar-change-adjunctions}
设 \(\varphi:R\to S\) 是保单位元的环同态，\(M\) 是左 \(R\) 模，\(N\) 是左 \(S\) 模。函子 \(\operatorname{Res}_{\varphi}:S\text{-}\mathbf{Mod}\to R\text{-}\mathbf{Mod}\) 以 \(r\cdot n=\varphi(r)n\) 限制标量。其左伴随为 \(S\otimes_R-\)，右伴随为
\[
\operatorname{Coind}_{\varphi}(M)
\coloneqq\operatorname{Hom}_R({}_RS,M),
\]
其中 \({}_RS\) 的左作用为 \(r\cdot s=\varphi(r)s\)，而这个 Hom 的左 \(S\) 作用为
\begin{equation}\label{b2:eq:coinduced-action}
(s f)(t)=f(ts).
\end{equation}
右伴随称为\term[yuyoudao]{余诱导}[coinduction]。具体地，有自然的群同构
\[
\begin{aligned}
\operatorname{Hom}_S(S\otimes_RM,N)
&\cong\operatorname{Hom}_R(M,\operatorname{Res}_{\varphi}N),\\
\operatorname{Hom}_S\bigl(N,\operatorname{Coind}_{\varphi}(M)\bigr)
&\cong\operatorname{Hom}_R(\operatorname{Res}_{\varphi}N,M).
\end{aligned}
\]
\end{proposition}
\begin{proof}
第一个同构把 \(h\) 送到 \(m\mapsto h(1\otimes m)\)，反向把 \(u:M\to\operatorname{Res}_{\varphi}N\) 送到
\[
s\otimes m\longmapsto s\cdot u(m).
\]
它平衡，因为 \(s\cdot\varphi(r)u(m)=s\cdot u(rm)\)，而且左 \(S\) 线性。两个公式互逆：一边使用 \(1u(m)=u(m)\)，另一边使用 \(h(s\otimes m)=s\cdot h(1\otimes m)\)。这也是\cref{b2:thm:hom-tensor-adjunction}取 \(P=S\) 的特例，其中 \(\operatorname{Hom}_S(S,N)\to N\) 的识别为在 \(1\) 处取值。

右伴随的标量作用同样来自预复合：令 \(R_s:S\to S\)、\(t\mapsto ts\)，则 \(sf=f\circ R_s\)。右乘 \(s\) 与左 \(R\) 作用可交换；具体地，若 \(f\) 左 \(R\) 线性，则
\[
(sf)\bigl(\varphi(r)t\bigr)=f\bigl(\varphi(r)ts\bigr)=r\cdot f(ts),
\]
故 \(sf\) 仍左 \(R\) 线性。又有
\(\bigl(s_1(s_2f)\bigr)(t)=f(ts_1s_2)=\bigl((s_1s_2)f\bigr)(t)\)，并有单位和分配律，所以确为左 \(S\) 模。

给定 \(v:N\to\operatorname{Coind}_{\varphi}(M)\)，将它送到 \(n\mapsto v(n)(1)\)。该映射左 \(R\) 线性，因为
\[
v\bigl(\varphi(r)n\bigr)(1)
=\bigl(\varphi(r)v(n)\bigr)(1)
=v(n)\bigl(\varphi(r)\bigr)=r\cdot v(n)(1).
\]
反向从左 \(R\) 线性的 \(u:\operatorname{Res}_{\varphi}N\to M\) 构造
\[
\widehat u(n)(s)=u(sn).
\]
对 \(s\) 的左 \(R\) 线性来自 \(u\bigl(\varphi(r)sn\bigr)=r\cdot u(sn)\)；对 \(n\) 的左 \(S\) 线性来自
\[
\widehat u(tn)(s)=u(stn)=\bigl(t\cdot\widehat u(n)\bigr)(s).
\]
在 \(1\) 处取值恢复 \(u(n)\)。反方向得到的函数在 \(s\) 处的值是
\(v(sn)(1)=\bigl(s\cdot v(n)\bigr)(1)=v(n)(s)\)，也恢复 \(v\)。各公式与 \(M\) 上的后复合、\(N\) 上的前复合相容，因为两条路径分别都为相同的取值与映射复合，故为自然同构。
\end{proof}

以商同态 \(R\to R/I\) 为例，限制标量的像正是被双边理想 \(I\) 零化的左模。扩张标量会把 \(IM\) 商掉，下一节将证明 \((R/I)\otimes_RM\cong M/IM\)。余诱导则保留子模
\[
\{\,m\in M\mid Im=0\,\}.
\]
确实，每个 \(f:R/I\to M\) 由 \(f(\bar1)=m\) 决定，关系 \(I\bar1=0\) 要求 \(Im=0\)；反过来这样的 \(m\) 唯一给出 \(f(\bar r)=rm\)。这是商模与子模两种不同的操作：对 \(M=\mathbb Z\)、\(R/I=\mathbb Z/n\mathbb Z\)、\(n>1\)，前者为 \(\mathbb Z/n\mathbb Z\)，后者为零。

% !TeX root = ../../Book2.tex
% 纲要标题：### 6.4 正合性、局部化与标量扩张
% 纲要位置：计划纲要.md，第 1262 行。

\section{正合性、局部化与标量扩张}
\label{b2:sec:0604}

给模添上一组关系，等于取一个商模。张量积会把这些关系送到新的模中，因而商模仍可通过张量计算；但它可能使原来非零的元素变成零，从而破坏单射。两种现象分别解释张量积为何右正合、为何未必正合。局部化与基域扩张也是改变标量的过程，本节还将用张量积描述它们。

% 纲要要点（供目录核对）：
%
% * 张量积的右正合性
% * 有限展示张量后的矩阵计算
% * 张量积不保持单射的实例
% * 局部化、商模与基域扩张中的计算
% * 局部化保持短正合列
%
% ---
%

\subsection{张量积的右正合性}
\label{b2:subsec:0604:01}

\begin{lemma}\label{b2:lem:tensor-cokernel}
设 \(R\) 是含幺环，\(M\) 为幺性右 \(R\) 模，\(N\) 为幺性左 \(R\) 模，\(K\) 为 \(N\) 的子模，\(\iota:K\to N\) 为包含映射。令
\[
H=\operatorname{im}(\operatorname{id}_M\otimes\iota)
\subseteq M\otimes_RN,
\]
即所有 \(m\otimes k\)（\(m\in M,k\in K\)）生成的子群，则自然映射诱导同构
\[
(M\otimes_RN)/H\longrightarrow M\otimes_R(N/K),
\quad m\otimes n+H\longmapsto m\otimes(n+K).
\]
\end{lemma}
\begin{proof}
商映射 \(q:N\to N/K\) 诱导 \(\operatorname{id}_M\otimes q\)，它将 \(H\) 送到零，故给出所述映射。为构造逆，考虑
\[
(m,n+K)\longmapsto m\otimes n+H.
\]
若以 \(n+k\) 代替 \(n\)，像之差为 \(m\otimes k+H=0\)，所以不依赖代表元。双可加性继承自原张量积，而 \(mr\otimes n=m\otimes rn\) 给出平衡性。故此式诱导反向同态。两复合在各自的纯张量或纯张量剩余类上为恒等，所以互逆。
\end{proof}

\begin{definition}\label{b2:def:right-exact-functor}
设 \(\mathcal C,\mathcal D\) 是两个模范畴，\(F:\mathcal C\to\mathcal D\) 是协变加性函子，即 \(F(f+g)=F(f)+F(g)\)。如果 \(\mathcal C\) 中每个正合列
\[
A\longrightarrow B\longrightarrow C\longrightarrow0
\]
在 \(F\) 下仍给出正合列
\[
F(A)\longrightarrow F(B)\longrightarrow F(C)\longrightarrow0,
\]
就称 \(F\) 为\term[youzhenghe-hanzi]{右正合函子}[right exact functor]。这里的模范畴可以分别是不同环的左模或右模范畴。
\end{definition}

\begin{theorem}[张量积的右正合性]\label{b2:thm:tensor-right-exact}
设 \(R\) 是含幺环，\(M\) 为幺性右 \(R\) 模，幺性左 \(R\) 模序列
\[
N'\xrightarrow{f}N\xrightarrow{g}N''\longrightarrow0
\]
正合，则交换群序列
\[
M\otimes_RN'\xrightarrow{\operatorname{id}\otimes f}
M\otimes_RN\xrightarrow{\operatorname{id}\otimes g}
M\otimes_RN''\longrightarrow0
\]
正合。对任意幺性左 \(R\) 模 \(L\)，若幺性右 \(R\) 模序列 \(Q'\xrightarrow{u}Q\xrightarrow{v}Q''\to0\) 正合，则
\[
Q'\otimes_RL\xrightarrow{u\otimes\operatorname{id}_L}
Q\otimes_RL\xrightarrow{v\otimes\operatorname{id}_L}
Q''\otimes_RL\longrightarrow0
\]
也正合。因此 \(M\otimes_R-:R\text{-}\mathbf{Mod}\to\mathbf{Ab}\) 和 \(-\otimes_RL:\mathbf{Mod}\text{-}R\to\mathbf{Ab}\) 都是右正合函子。

若 \(S,T\) 是含幺环，\(M\) 为 \((S,R)\) 双模，\(N',N,N''\) 为 \((R,T)\) 双模且 \(f,g\) 为双模同态，则第一个序列在 \((S,T)\) 双模范畴中正合。若 \(Q',Q,Q''\) 为 \((S,R)\) 双模且 \(u,v\) 为双模同态，\(L\) 为 \((R,T)\) 双模，则第二个序列也在 \((S,T)\) 双模范畴中正合。
\end{theorem}
\begin{proof}
令 \(K=\operatorname{im}f=\ker g\)。由于 \(g\) 满射，\cref{b2:thm:module-first-isomorphism}给出 \(N/K\cong N''\)。\cref{b2:lem:tensor-cokernel}说明 \(\operatorname{id}\otimes g\) 的核恰为所有 \(m\otimes k\)、\(k\in K\) 生成的子群。对每个这样的 \(k\)，存在 \(n'\in N'\) 使 \(f(n')=k\)，故
\[
m\otimes k=(\operatorname{id}\otimes f)(m\otimes n').
\]
反过来，\(\operatorname{id}\otimes f\) 的每个纯张量像均为这种形式，于是其像恰为上述核。又因 \(g\) 满射，每个 \(m\otimes n''\) 都有原像 \(m\otimes n\)，纯张量生成目标，所以 \(\operatorname{id}\otimes g\) 满射。

第一变量的结论可将已证结论用于 \(R^{\mathrm{op}}\)，再用已证明的交换因子同构 \eqref{b2:eq:tensor-opposite-swap}转换回来：它将 \(Q'\otimes_RL\to Q\otimes_RL\) 对应于 \(L\otimes_{R^{\mathrm{op}}}Q'\to L\otimes_{R^{\mathrm{op}}}Q\)，因此第二变量的结论给出原来第一变量的正合性。\cref{b2:prop:tensor-maps}已说明这两个张量函子保持同态的加法。外侧作用分别由 \(s(x\otimes y)=(sx)\otimes y\) 与 \((x\otimes y)t=x\otimes(yt)\) 给出；相应线性条件保证张量后的箭头保留作用。模的核与像在底交换群上计算，故上述序列在所述双模范畴中正合。
\end{proof}

“右正合”中的“右”指序列末端以零结束，不指右模。在定理的记号下，\(\operatorname{coker}f=N/\operatorname{im}f\)，因而上述结论也写成
\[
M\otimes_R\operatorname{coker}f
\cong\operatorname{coker}(\operatorname{id}_M\otimes f).
\]
张量积保留的是由关系生成的商；它没有断言 \(\operatorname{id}\otimes f\) 单射，即使原来的 \(f\) 已经单射。只有左侧或右侧外侧作用时，可分别在定理中取 \(T=\mathbb Z\) 或 \(S=\mathbb Z\)，利用每个交换群的自然整数作用，得到左模或右模范畴中的相应结论。

\ProseParagraphBreak
对环 \(R\) 的双边理想 \(I\)，若要让左 \(R\) 模 \(N\) 的作用经过商环 \(R/I\)，就必须让 \(I\) 的每个元素作用为零。在张量积中，平衡关系给出 \(\bar1\otimes in=\bar i\otimes n=0\)，因此 \(IN\) 中的有限和都被消去。下面的同构说明，无须再加其他关系；右模的情形相同。

\begin{proposition}[商环上的张量计算]\label{b2:prop:tensor-quotient}
设 \(R\) 是含幺环，\(I\) 为 \(R\) 的双边理想，\(N\) 为幺性左 \(R\) 模，\(M\) 为幺性右 \(R\) 模。存在自然同构
\[
(R/I)\otimes_RN\cong N/IN,
\qquad M\otimes_R(R/I)\cong M/MI,
\]
具体公式分别为 \(\bar r\otimes n\mapsto rn+IN\) 与 \(m\otimes\bar r\mapsto mr+MI\)。所得对象分别带左、右 \(R/I\) 模结构。
\end{proposition}
\begin{proof}
第一式的公式不依赖 \(r\) 的代表元，因为 \(i n\in IN\)；它双可加且平衡，所以诱导同态。反向公式为 \(n+IN\mapsto\bar1\otimes n\)。它良定义：\(IN\) 中的元素是有限和 \(\sum_j i_j n_j\)，而
\[
\bar1\otimes\sum_j i_j n_j
=\sum_j\bar i_j\otimes n_j=0.
\]
两个公式互逆，因为 \(\bar1\otimes rn=\bar r\otimes n\)，且 \(1n=n\)。左 \(R/I\) 线性由公式直接给出。第二式用 \(m+MI\mapsto m\otimes\bar1\) 作逆，\(m_j i_j\otimes\bar1=m_j\otimes\bar i_j=0\) 核对关系；其余复合与右线性同样由给出的公式逐项得到。
\end{proof}

对交换环 \(R\) 的两个理想 \(I,J\)，于是有
\[
(R/I)\otimes_R(R/J)\cong R/(I+J).
\]
确实，上一命题先给出 \((R/J)/I(R/J)\)，而 \(I(R/J)=(I+J)/J\)，再用商模的迭代取商得到结果。特别地，正整数 \(m,n\) 满足
\[
(\mathbb Z/m\mathbb Z)\otimes_{\mathbb Z}(\mathbb Z/n\mathbb Z)
\cong\mathbb Z/\gcd(m,n)\mathbb Z,
\]
因为整数理想之和 \(m\mathbb Z+n\mathbb Z=\gcd(m,n)\mathbb Z\)。
取 \(m=2,n=3\)，便得到本章开头的零张量积，这是互素整数理想之和为整个 \(\mathbb Z\) 的结果。

\subsection{张量积不保持单射的例子}
\label{b2:subsec:0604:02}

取整数 \(n>1\)，序列
\[
0\longrightarrow\mathbb Z\xrightarrow{\times n}\mathbb Z
\longrightarrow\mathbb Z/n\mathbb Z\longrightarrow0
\]
正合。与 \(\mathbb Z/n\mathbb Z\) 张量，利用\cref{b2:thm:tensor-associativity-unit}的单位同构及\cref{b2:prop:tensor-quotient}，得到
\[
\mathbb Z/n\mathbb Z\xrightarrow{0}\mathbb Z/n\mathbb Z
\xrightarrow{\operatorname{id}}\mathbb Z/n\mathbb Z\longrightarrow0.
\]
第一个映射为零，因为 \(n a\otimes\bar b=a\otimes n\bar b=0\)；原先的“乘 \(n\)”在张量后仍是乘同一个整数，但此时 \(n\) 已经作用为零。它的来源非零，所以不再单射。第二个映射在上述识别下为恒等，故序列末端仍正合，完全符合右正合定理。

\ProseParagraphBreak
即使张量后单射仍然成立，也不能只凭原映射是包含就不加检查。一个足够条件是原包含有缩回。若 \(i:N'\to N\)、\(r:N\to N'\) 满足 \(ri=\operatorname{id}\)，则
\[
(\operatorname{id}_M\otimes r)\circ(\operatorname{id}_M\otimes i)
=\operatorname{id}_{M\otimes_RN'}
\]
由\cref{b2:prop:tensor-maps}成立，因此张量后的 \(i\) 仍单射，且仍有缩回。结合右正合性，张量积保持每个分裂短正合列。

\ProseParagraphBreak
非分裂时，单射是否保留取决于用来张量的模。下一章的平坦模就是把“对所有单射都保留单射”作为条件。本节的 \(\mathbb Z/n\mathbb Z\) 展示了该条件并非自动满足，不能用“张量积类似向量空间的乘积”代替核的实际计算。

\subsection{局部化与基域扩张}
\label{b2:subsec:0604:03}

商模公式通过让理想中的元素作用为零来改变标量。局部化则让指定的标量成为可逆元；以下只讨论交换环，此时无须区分标量从左还是从右作用。

\ProseParagraphBreak
把分母变成单位，也会使被分母零化的元素变成零。例如，对 \(R=\mathbb Z\)、\(\Sigma=\{1,2,2^2,\ldots\}\)、\(N=\mathbb Z/2\mathbb Z\)，局部化后 \(2\) 可逆，而 \(2\bar1=0\)，所以 \(\bar1\) 的像必须为零。若只按交叉相乘相等识别分式，\((\bar1,1)\) 与 \((0,1)\) 却不能等价，因为原模中 \(\bar1\ne0\)。因此必须允许交叉乘积之差再被某个分母零化。

\begin{definition}\label{b2:def:localized-module}
设 \(R\) 是交换含幺环，\(\Sigma\subseteq R\) 是含 \(1\) 的乘法封闭集，\(N\) 是幺性 \(R\) 模。在 \(N\times\Sigma\) 上规定关系
\[
(n,s)\sim(n',s')\quad\Longleftrightarrow\quad
\text{存在}\;u\in\Sigma\;\text{使}\;u(s'n-sn')=0.
\]
将 \((n,s)\) 的等价类记为 \(n/s\)。在全体等价类上定义以下加法与环局部化 \(\Sigma^{-1}R\) 的标量作用：
\[
\frac ns+\frac{n'}t=\frac{tn+sn'}{st},\qquad
\frac av\frac ns=\frac{an}{vs}.
\]
所得 \(\Sigma^{-1}R\) 模称为 \(N\) 在 \(\Sigma\) 处的\term[mojubu-hua]{局部化}[localization of a module] \(\Sigma^{-1}N\)。其中 \(\Sigma^{-1}R\) 是第一卷定理13.2.2所构造的交换环局部化。
\end{definition}
\symidx{\(\Sigma^{-1}N\)：模 \(N\) 的局部化}

先验证关系确实是等价关系。反身性与对称性直接成立。若 \(u(tn-sn')=0\)、\(v(wn'-tn'')=0\)，则把第一式乘 \(vw\)，第二式乘 \(us\) 后相加，得到
\(uvt(wn-sn'')=0\)。因 \(uvt\in\Sigma\)，传递性成立。

\ProseParagraphBreak
关系中的额外因子也使运算能容许被分母零化的误差。若 \(u(s'n-sn_1)=0\)、\(v(t'n'-tn_1')=0\)，则两种加法结果的交叉相乘之差为
\[
t t'(s'n-sn_1)+s s'(t'n'-tn_1'),
\]
它被 \(uv\) 零化。对标量乘法，若 \(w(v'a-va')=0\)，模分式的关系如前，则两结果的交叉相乘之差可以写成
\[
v'a(s'n-sn_1)+s(v'a-va')n_1,
\]
它被 \(uw\) 零化，故两项运算都不依赖代表元。三个模分式通分到分母的乘积后，两种相加次序有同一分子，因而加法结合律成立；零元是 \(0/1\)，负元是 \((-n)/s\)。标量作用的结合律则由
\[
\frac av\left(\frac bw\frac ns\right)
=\frac{abn}{vws}
=\left(\frac av\frac bw\right)\frac ns
\]
给出，单位作用及两条分配律同样在通分后化为 \(R\) 与 \(N\) 的对应公理。这就得到所述局部化模。

\ProseParagraphBreak
一个分式何时为零，只需在等价关系中与 \((0,1)\) 比较：
\begin{equation}\label{b2:eq:localization-zero}
\frac ns=0\quad\Longleftrightarrow\quad
\text{存在}\;u\in\Sigma\;\text{使}\;un=0.
\end{equation}
因此典范映射 \(N\to\Sigma^{-1}N\)、\(n\mapsto n/1\) 的核，恰为被某个分母零化的元素；局部化没有再消去其他元素。

\begin{proposition}[局部化的张量描述]\label{b2:prop:localization-tensor}
设 \(R\) 是交换含幺环，\(\Sigma\subseteq R\) 是含 \(1\) 的乘法封闭集，\(N\) 是幺性 \(R\) 模，则存在自然的 \(\Sigma^{-1}R\) 模同构
\[
(\Sigma^{-1}R)\otimes_RN\longrightarrow\Sigma^{-1}N,
\qquad \frac as\otimes n\longmapsto\frac{an}s.
\]
其逆映射为 \(n/s\mapsto(1/s)\otimes n\)。
\end{proposition}
\begin{proof}
正向公式对分式 \(a/s\) 的代表元良定义：若 \(u(ta-sb)=0\)，则 \(u(tan-sbn)=0\)。分式加法与模加法给出双可加性，\((ar)n=a(rn)\) 给出 \(R\) 平衡性，故诱导同态，并且保留 \(\Sigma^{-1}R\) 作用。

若 \(n/s=n'/t\)，取 \(u\in\Sigma\) 使 \(u(tn-sn')=0\)。在张量积中，
\[
\frac1s\otimes n-\frac1t\otimes n'
=\frac1{ust}\otimes u(tn-sn')=0.
\]
于是逆映射不依赖代表元。通分公式说明它保持加法，标量作用则由
\[
\frac1{vs}\otimes an=\frac a{vs}\otimes n
=\frac av\left(\frac1s\otimes n\right)
\]
核对。两复合分别恢复 \(n/s\) 与 \((a/s)\otimes n\)，所以互逆。对模同态 \(f:N\to N'\)，两种构造中对应的映射分别为 \((a/s)\otimes n\mapsto(a/s)\otimes f(n)\) 与 \(n/s\mapsto f(n)/s\)，故同构自然。
\end{proof}

\begin{proposition}\label{b2:prop:localization-exact}
设 \(R\) 是交换含幺环，\(\Sigma\subseteq R\) 是含 \(1\) 的乘法封闭集，幺性 \(R\) 模的短正合列为
\[
0\longrightarrow N'\xrightarrow{f}N\xrightarrow{g}N''\longrightarrow0.
\]
则 \(\Sigma^{-1}R\) 模序列
\[
0\longrightarrow\Sigma^{-1}N'
\xrightarrow{\Sigma^{-1}f}\Sigma^{-1}N
\xrightarrow{\Sigma^{-1}g}\Sigma^{-1}N''\longrightarrow0
\]
仍正合，其中局部化同态的公式为 \((\Sigma^{-1}f)(n'/s)=f(n')/s\)、\((\Sigma^{-1}g)(n/s)=g(n)/s\)。
\end{proposition}
\begin{proof}
对任意 \(R\) 模同态 \(h\)，若 \(u(tn-sn_1)=0\)，则 \(u(th(n)-sh(n_1))=h(u(tn-sn_1))=0\)，所以 \(n/s\mapsto h(n)/s\) 良定义；分式运算说明它保留加法与 \(\Sigma^{-1}R\) 作用。由\cref{b2:prop:localization-tensor}的自然性，序列中这两个同态对应于 \(\operatorname{id}_{\Sigma^{-1}R}\otimes f\)、\(\operatorname{id}_{\Sigma^{-1}R}\otimes g\)。\cref{b2:thm:tensor-right-exact}因此保证末端正合。

还须证明 \(\Sigma^{-1}f\) 单射。若 \(f(n')/s=0\)，由\eqref{b2:eq:localization-zero}，存在 \(u\in\Sigma\) 使 \(uf(n')=f(un')=0\)。原映射 \(f\) 单射，故 \(un'=0\)，再由同一零元判据得 \(n'/s=0\)。于是整个局部化序列短正合。
\end{proof}

局部化不仅保留由关系得到的商，还保留单射；下一章将把这项性质表述为 \(\Sigma^{-1}R\) 的平坦性。

\ProseParagraphBreak
例如，\(\Sigma=\{1,p,p^2,\ldots\}\subseteq\mathbb Z\) 时，\(\Sigma^{-1}\mathbb Z=\mathbb Z[1/p]\)。若 \(N=\mathbb Z/p^a\mathbb Z\)，每个元素被 \(p^a\) 零化，由\eqref{b2:eq:localization-zero}可知局部化为零。若 \(\gcd(p,n)=1\)，乘以 \(p\) 已经是 \(\mathbb Z/n\mathbb Z\) 的自同构，分式 \(\bar a/p^j\) 对应 \(p^{-j}\bar a\)，所以此时局部化仍为 \(\mathbb Z/n\mathbb Z\)。

\ProseParagraphBreak
取所有非零整数作分母，得到 \(\mathbb Q\otimes_{\mathbb Z}N\)。由零元判据，典范映射 \(N\to\mathbb Q\otimes_{\mathbb Z}N\)、\(n\mapsto1\otimes n\) 的核恰为挠子模；但有理化并不只是商去挠子模，还允许剩余元素带任意非零整数分母，形成一个 \(\mathbb Q\) 向量空间。例如 \(N=\mathbb Z\) 本来无挠，而 \(\mathbb Q\otimes_{\mathbb Z}\mathbb Z\cong\mathbb Q\)，并不等于 \(\mathbb Z\)。

\ProseParagraphBreak
局部化把环 \(R\) 换成 \(\Sigma^{-1}R\)。另一种标量变换是把域 \(k\) 扩张到更大的域 \(K\)。取域扩张 \(k\subseteq K\) 与 \(k\) 向量空间 \(V\)，则 \(K\otimes_kV\) 是 \(K\) 向量空间，称为 \(V\) 的\term[jiyukuozhang]{基域扩张}[extension of the ground field]。若 \((v_i)_{i\in I}\) 是 \(V\) 的基，则
\[
(1\otimes v_i)_{i\in I}
\]
是 \(K\otimes_kV\) 的 \(K\) 基。事实上，先用交换因子同构，再用\cref{b2:prop:tensor-free-coordinate}，得到坐标同构 \(K\otimes_kV\cong K^{(I)}\)，其公式正是把 \(\sum_i a_i\otimes v_i\) 送到 \((a_i)\)。所以 \(v\mapsto1\otimes v\) 单射，且
\(\dim_K(K\otimes_kV)=\dim_kV\)。这里比较的是两个不同基域上的维数。

\ProseParagraphBreak
线性映射 \(T:V\to W\) 扩张为 \(\operatorname{id}_K\otimes T\)，它在扩张后的基下保留原矩阵。例如，实线性算子
\[
T=\begin{pmatrix}0&-1\\1&0\end{pmatrix}
\]
没有实特征值；在 \(\mathbb C\otimes_{\mathbb R}\mathbb R^2\cong\mathbb C^2\) 上，同一个矩阵有特征向量 \((1,-i)\)、\((1,i)\)，特征值分别为 \(i,-i\)。两个向量组成的矩阵行列式为 \(2i\ne0\)，故形成复基，使扩张算子对角化。基域扩张保留线性映射的公式，却可以改变多项式的分解情况，进而改变可用的标准形。

\ProseParagraphBreak
基域扩张保留自由模的基；对于有限展示模，则可以直接把关系矩阵的系数作用于新的系数模。这也给出了不要求新系数模平坦的张量计算。

\begin{corollary}\label{b2:cor:tensor-finite-presentation}
设 \(R\) 是交换含幺环，\(M,N\) 是幺性 \(R\) 模，且 \(M\) 有有限展示
\[
R^m\xrightarrow{A}R^n\xrightarrow{\pi}M\longrightarrow0,
\qquad A=(a_{ij})\in M_{n\times m}(R),
\]
其中 \(m,n\) 为非负整数，向量均按列书写。定义 \(R\) 模同态
\[
A_N:N^m\longrightarrow N^n,\qquad
(z_j)_{j=1}^m\longmapsto
\left(\sum_{j=1}^m a_{ij}z_j\right)_{i=1}^n.
\]
若 \(e_1,\ldots,e_n\) 为 \(R^n\) 的标准基，则序列
\[
N^m\xrightarrow{A_N}N^n\longrightarrow M\otimes_RN\longrightarrow0,
\qquad (z_i)\longmapsto\sum_{i=1}^n\pi(e_i)\otimes z_i,
\]
正合，并给出自然的 \(R\) 模同构
\[
N^n/\operatorname{im}A_N\longrightarrow M\otimes_RN,
\qquad (z_i)+\operatorname{im}A_N\longmapsto
\sum_{i=1}^n\pi(e_i)\otimes z_i.
\]
零次直和为零模，空和按零理解。
\end{corollary}
\begin{proof}
对展示在第一变量中张量 \(N\)，由\cref{b2:thm:tensor-right-exact}得到
\[
R^m\otimes_RN\xrightarrow{A\otimes\operatorname{id}_N}
R^n\otimes_RN\xrightarrow{\pi\otimes\operatorname{id}_N}
M\otimes_RN\longrightarrow0.
\]
\cref{b2:prop:tensor-free-coordinate}把前两项分别识别为 \(N^m,N^n\)，其逆公式为 \((z_i)\mapsto\sum_i e_i\otimes z_i\)；交换性保证这些坐标同构保留 \(R\) 作用。若 \(\widetilde e_j\) 为 \(R^m\) 的第 \(j\) 个标准基向量，则
\[
(A\otimes\operatorname{id}_N)(\widetilde e_j\otimes z)
=\sum_{i=1}^n e_i a_{ij}\otimes z
=\sum_{i=1}^n e_i\otimes a_{ij}z.
\]
故第一个映射在坐标下恰为 \(A_N\)，第二个映射恰为所述有限求和。正合性说明后者满射且核为 \(\operatorname{im}A_N\)，再由\cref{b2:thm:module-first-isomorphism}得到商模同构。对任意 \(R\) 模同态 \(h:N\to N'\)，逐坐标施加 \(h\) 与 \(A_N,A_{N'}\) 相容，并把上述有限和送到 \(\sum_i\pi(e_i)\otimes h(z_i)\)，这证明自然性。
\end{proof}


\begin{chaptersummary}
张量积由纯张量生成，但生成元之间受双可加性与平衡关系约束。从它出发定义映射，最可靠的方法是先在两变量上写公式，再核对这些关系。一般环上的张量积首先是交换群；双模提供的外侧作用才使它成为新的模。

\ProseParagraphBreak
结合同构与单位同构使双模张量可以连续进行，反环则记录右作用对应的乘法次序。Hom–张量伴随来自同一个双变量公式的两种读法，其自然性落实为前复合与后复合的明确等式。限制标量同时有扩张标量和余诱导两个伴随，二者通常产生不同对象。

\ProseParagraphBreak
张量积保留余核，因而右正合；单射却可能在张量后变成零映射。商环张量使理想作用归零，局部化使分母可逆，基域扩张则保留基的大小和线性映射的矩阵。下一章将据此研究哪些模能使张量积始终保持正合，以及它们与自由模、投射模之间的关系。
\end{chaptersummary}

\BookExercises

\begin{exercise}\label{b2:ex:06-balanced-cyclic}
设 \(m,n>0\)，\(A\) 为交换群。证明从 \((\mathbb Z/m\mathbb Z)\times(\mathbb Z/n\mathbb Z)\) 到 \(A\) 的平衡映射由 \(x=b(\bar1,\bar1)\) 唯一决定，并求出 \(x\) 必须满足的条件。由此重新求出这两个循环群的张量积。
\end{exercise}
\begin{solution}
双可加性迫使 \(b(\bar a,\bar c)=acx\)。代表元 \(a\) 加上 \(m\)、\(c\) 加上 \(n\) 时不改变结果，当且仅当 \(mx=nx=0\)。令 \(d=\gcd(m,n)\)，由 Bézout 等式，这等价于 \(dx=0\)：一方向将 \(d\) 写成 \(m,n\) 的整数线性组合，另一方向用 \(d\mid m,n\)。对满足条件的 \(x\)，公式确实双可加且平衡。取 \(x=\bar1\in\mathbb Z/d\mathbb Z\)，得到平衡映射 \((\bar a,\bar c)\mapsto\overline{ac}\)。每个上述 \(b\) 都唯一地经由它分解，因为从 \(\mathbb Z/d\mathbb Z\) 出发的同态由满足 \(dx=0\) 的 \(\bar1\) 的像唯一决定。由\cref{b2:thm:tensor-universal}，这给出所求同构，\(\bar1\otimes\bar1\) 对应 \(\bar1\)。
\end{solution}

\begin{exercise}\label{b2:ex:06-nonzero-pure}
设 \(V,W\) 是域 \(k\) 上的向量空间。证明 \(v\otimes w=0\) 当且仅当 \(v=0\) 或 \(w=0\)。给出一般环上两因子均非零、纯张量却为零的例子。
\end{exercise}
\begin{solution}
任一因子为零时由张量关系得到结论。若 \(v,w\ne0\)，分别将它们扩充为基，定义线性泛函 \(f:V\to k\)、\(g:W\to k\)，使 \(f(v)=g(w)=1\)。平衡映射 \((x,y)\mapsto f(x)g(y)\) 诱导线性映射 \(V\otimes_kW\to k\)，将 \(v\otimes w\) 送到 \(1\)，故该张量非零。一般环上取 \(R=\mathbb Z\)、\(V=\mathbb Z/2\mathbb Z\)、\(W=\mathbb Z/3\mathbb Z\)；两个 \(\bar1\) 都非零，而其张量同时被 \(2,3\) 零化，所以为零。
\end{solution}

\begin{exercise}\label{b2:ex:06-matrix-rank-one}
以 \(k^a,k^b\) 的标准基写张量的坐标矩阵。证明非零张量为纯张量，当且仅当其坐标矩阵秩为 \(1\)。据此判断 \(e_1\otimes f_1+e_2\otimes f_2\) 是否为纯张量。
\end{exercise}
\begin{solution}
若 \(v=\sum_i a_ie_i\)、\(w=\sum_j b_jf_j\)，则 \(v\otimes w\) 的矩阵为 \((a_ib_j)\)。非零时它的每一列都是非零向量 \((a_i)\) 的倍数，至少一列非零，所以秩为 \(1\)。反过来，秩为 \(1\) 的矩阵可取一条非零列 \(a\)，将其他列写成 \(b_j a\)，从而矩阵为 \((a_i b_j)\)，对应 \(v\otimes w\)。这里使用\cref{b2:prop:tensor-free-coordinate}的坐标唯一性。题末张量的矩阵左上有 \(2\times2\) 单位子矩阵，其余条目为零，所以秩为 \(2\)，不是纯张量。
\end{solution}

\begin{exercise}\label{b2:ex:06-direct-sum}
计算
\[
\left(\mathbb Z^2\oplus\mathbb Z/12\mathbb Z\oplus\mathbb Z/18\mathbb Z\right)
\otimes_{\mathbb Z}\mathbb Z/8\mathbb Z.
\]
写出各循环因子的阶，并说明结果为什么不是简单地把原来的每个整数模数都改成 \(8\)。
\end{exercise}
\begin{solution}
由\cref{b2:prop:tensor-direct-sums}分开四个直和分量。两个自由分量分别给出 \(\mathbb Z/8\mathbb Z\)；循环分量由\cref{b2:prop:tensor-quotient}给出模数 \(\gcd(12,8)=4\) 与 \(\gcd(18,8)=2\)。因此结果为
\[
(\mathbb Z/8\mathbb Z)^2\oplus\mathbb Z/4\mathbb Z\oplus\mathbb Z/2\mathbb Z.
\]
原来的关系 \(12x=0\)、\(18y=0\) 仍须保留，同时增添 \(8x=8y=0\)，共同得到更小的阶。扩张标量并不删除原有关系。
\end{solution}

\begin{exercise}\label{b2:ex:06-matrix-bimodule}
令 \(R=M_n(k)\)，\(n\ge1\)，将行向量空间 \(U=k^{1\times n}\) 看成右 \(R\) 模，将列向量空间 \(V=k^{n\times1}\) 看成左 \(R\) 模。证明矩阵乘法诱导 \(U\otimes_RV\cong k\)。再证明外积给出 \((R,R)\) 双模同构 \(V\otimes_kU\cong R\)。
\end{exercise}
\begin{solution}
设 \(u_i,v_j\) 为行、列标准基。乘法 \((u,v)\mapsto uv\in k\) 平衡，所以诱导映射。由 \(u_iE_{ij}=u_j\)，有
\[
u_j\otimes v_\ell=u_iE_{ij}\otimes v_\ell
=u_i\otimes E_{ij}v_\ell=\delta_{j\ell}u_i\otimes v_i.
\]
因此所有非对角纯基张量为零，所有对角纯基张量相等。其共同像为 \(1\)，故 \(c\mapsto c u_1\otimes v_1\) 与乘法映射互逆。

第二个映射将 \(v_i\otimes u_j\) 送到矩阵单位 \(E_{ij}\)。由\cref{b2:prop:tensor-free-coordinate}，这些纯基张量是 \(k\) 基，矩阵单位也是 \(R\) 的 \(k\) 基，所以映射为同构。等式 \((Av)u=A(vu)\)、\(v(uB)=(vu)B\) 表明它保留两侧的 \(R\) 作用。

两个同构保留的作用不同：\(U\otimes_RV\) 是 \(k\) 上的对象，\(V\otimes_kU\) 则仍有两侧的矩阵作用。第16章的 Morita 理论将用这种成对的双模及其张量同构比较模范畴。
\end{solution}

\begin{exercise}\label{b2:ex:06-unit-naturality}
对右 \(R\) 模同态 \(f:M\to M'\)，证明单位同构满足
\(\rho_{M'}\circ(f\otimes\operatorname{id}_R)=f\circ\rho_M\)。若 \(P\) 为 \((R,S)\) 双模，证明将 \((M\otimes_RR)\otimes_RP\) 识别为 \(M\otimes_RP\) 的两条路径相同：一条先用 \(\rho_M\)，另一条先用结合同构再用 \(R\otimes_RP\cong P\)。
\end{exercise}
\begin{solution}
第一式两边在 \(m\otimes r\) 上分别为 \(f(m)r\) 与 \(f(mr)\)，由右线性相等。第二问对生成元 \((m\otimes r)\otimes p\)，第一条路径给出 \(mr\otimes p\)，第二条给出 \(m\otimes rp\)；由平衡关系相同。生成元上的相等延伸为整个群上的相等，说明这些识别与同态及结合规则相容。
\end{solution}

\begin{exercise}\label{b2:ex:06-opposite-ring}
对任意环 \(R\)，证明 \(\operatorname{End}_{R^{\mathrm{op}}}(R_R)\cong R\)。将它与 \(\operatorname{End}_R({}_RR)\cong R^{\mathrm{op}}\) 比较，明确两种情形中元素如何作用，以及复合的次序。
\end{exercise}
\begin{solution}
右线性映射满足 \(f(r)=f(1)r\)，故由 \(a=f(1)\) 决定，并且等于左乘 \(L_a\)。复合满足 \(L_aL_b(r)=a(br)=(ab)r\)，所以 \(a\mapsto L_a\) 保持乘法、加法与单位，且双射。左线性情形则为右乘 \(D_a(r)=ra\)，有 \(D_aD_b(r)=(rb)a=r(ba)\)。因此前一种给原环，后一种给反环，差异来自左乘与右乘的复合次序。
\end{solution}

\begin{exercise}\label{b2:ex:06-adjunction-evaluation}
设 \({}_SP_R\) 为双模、\(N\) 为左 \(S\) 模。直接证明求值公式
\(p\otimes f\mapsto f(p)\) 对 \(R\) 平衡且左 \(S\) 线性，并求出它在 Hom–张量伴随下所对应的映射。
\end{exercise}
\begin{solution}
两变量可加，且 \(f(pr)=(rf)(p)\)，所以 \((pr,f)\) 与 \((p,rf)\) 的像相同。由\cref{b2:thm:tensor-universal}，求值诱导群同态。又 \(f(sp)=s\cdot f(p)\)，所以它左 \(S\) 线性。采用\cref{b2:thm:hom-tensor-adjunction}并取 \(M=\operatorname{Hom}_S(P,N)\)，相应映射将 \(f\) 送到函数 \(p\mapsto f(p)\)，正是 \(\operatorname{Hom}_S(P,N)\) 的恒等映射。
\end{solution}

\begin{exercise}\label{b2:ex:06-hom-tensor-finite-free}
设 \(k\) 为域，\(P=\bigoplus_{n\ge0}k e_n\)，\(P^*=\operatorname{Hom}_k(P,k)\)。证明自然映射
\[
\Theta:P^*\otimes_kP\longrightarrow\operatorname{End}_k(P),
\qquad \lambda\otimes v\longmapsto\bigl(p\mapsto\lambda(p)v\bigr)
\]
单射，且其像恰为像空间有限维的自同态。由此说明恒等自同态不在像中，并与\cref{b2:prop:infinite-free-hom-tensor}中的共同分母限制比较：这里没有分母，有限张量和仍带来了什么限制？
\end{exercise}
\begin{solution}
有限和 \(\sum_{j=1}^r\lambda_j\otimes v_j\) 所定义的映射，其像包含于有限维空间 \(\operatorname{span}_k(v_1,\ldots,v_r)\)。反过来，若 \(h(P)\) 有有限基 \(v_1,\ldots,v_r\)，将 \(h(p)\) 在此基中的坐标记为 \(\lambda_j(p)\)。坐标泛函与 \(h\) 的复合均为线性泛函，且 \(h(p)=\sum_j\lambda_j(p)v_j\)，故 \(h=\Theta(\sum_j\lambda_j\otimes v_j)\)。

任一张量均可通过重组选取有限维的第二因子空间的一组基，写成 \(\sum_{j=1}^r\lambda_j\otimes v_j\)，其中 \(v_j\) 线性无关。若其像为零，则对每个 \(p\) 有 \(\sum_j\lambda_j(p)v_j=0\)，从而所有 \(\lambda_j(p)=0\)，即所有 \(\lambda_j=0\)。因此 \(\Theta\) 单射。

恒等映射的像是无限维的 \(P\)，故不在像中。共同分母反例限制的是所有取值能否由同一批分母表示；这里限制的是所有像能否落在同一个有限维子空间中。两者都来自一个张量只能使用有限项。
\end{solution}

\begin{exercise}\label{b2:ex:06-coinduction-quotient}
取 \(R=\mathbb Z\)、\(S=\mathbb Z/6\mathbb Z\)、\(M=\mathbb Z/12\mathbb Z\)。分别计算扩张标量 \(S\otimes_RM\) 与余诱导 \(\operatorname{Hom}_R(S,M)\)，明确它们的 \(S\) 作用。二者是否同构？这是否意味着扩张标量与余诱导是同一个函子？
\end{exercise}
\begin{solution}
扩张标量为 \(M/6M\)，由 \(\bar1\) 的剩余类生成，阶为 \(6\)，故同构于 \(S\)。余诱导由 \(f(\bar1)\) 决定，这个像必须被 \(6\) 零化；在 \(\mathbb Z/12\mathbb Z\) 中恰是偶数剩余类，组成由 \(\bar2\) 生成的阶 \(6\) 子群。两者的 \(S\) 作用均为相应整数倍，因此这里确实同构。可是换成 \(M=\mathbb Z\) 时，扩张标量为 \(S\)，余诱导为 \(0\)。一个对象处的同构既不说明两个函子相等，也不说明它们自然同构。
\end{solution}

\begin{exercise}\label{b2:ex:06-coinduction-rational}
对 \(\mathbb Z\to\mathbb Q\) 计算 \(\mathbb Q\otimes_{\mathbb Z}\mathbb Z\) 与 \(\operatorname{Hom}_{\mathbb Z}(\mathbb Q,\mathbb Z)\)，并据此说明限制标量的两个伴随可能相差很大。
\end{exercise}
\begin{solution}
第一项由单位同构等于 \(\mathbb Q\)。对任意加法同态 \(f:\mathbb Q\to\mathbb Z\)、\(q\in\mathbb Q\) 和正整数 \(n\)，有 \(f(q)=n\cdot f(q/n)\)，故整数 \(f(q)\) 被所有正整数整除，只能为零。于是整个 Hom 群为零。其左 \(\mathbb Q\) 作用也是零模的唯一作用。
\end{solution}

\begin{exercise}\label{b2:ex:06-tensor-presentation}
令 \(M\) 为整数矩阵 \(A=\begin{pmatrix}2&0\\1&3\end{pmatrix}\) 的余核。对 \(N=\mathbb Z/3\mathbb Z,\mathbb Z/4\mathbb Z,\mathbb Q\)，分别计算 \(M\otimes_{\mathbb Z}N\)，并在 \(N^2\) 上给出明确的商映射。若 \(x,y\) 是 \(M\) 中两个标准基向量的像，写出 \(x\otimes\bar1\)、\(y\otimes\bar1\) 在前两个计算结果中的坐标，解释两条原关系如何保留。
\end{exercise}
\begin{solution}
由\cref{b2:cor:tensor-finite-presentation}，结果是 \(N^2/A_NN^2\)。当 \(N=\mathbb Z/3\mathbb Z\) 时，矩阵为 \(\begin{pmatrix}2&0\\1&0\end{pmatrix}\)，像为 \(\langle(2,1)\rangle\)。商映射 \((a,b)\mapsto a+b\) 到 \(\mathbb Z/3\mathbb Z\) 的核恰为此像，故张量积为 \(\mathbb Z/3\mathbb Z\)，两个指定张量都对应 \(\bar1\)。

当 \(N=\mathbb Z/4\mathbb Z\) 时，第二列 \((0,3)\) 生成整个第二坐标，第一列再使第一坐标模去 \(2N\)。因此像为 \(2N\times N\)，商映射可取 \((a,b)\mapsto a\bmod2\)，结果为 \(\mathbb Z/2\mathbb Z\)。此时 \(x\otimes\bar1\) 对应 \(\bar1\)，\(y\otimes\bar1\) 对应零。

在 \(\mathbb Q\) 上，矩阵行列式为 \(6\ne0\)，故可逆，余核为零，商映射为零映射。在三个情形中，原关系 \(2x+y=0\)、\(3y=0\) 都成立；改变系数模使这些关系与新标量条件共同作用，并未删除原关系。
\end{solution}

\begin{exercise}\label{b2:ex:06-dual-number-kernel}
令 \(R=k[\varepsilon]/(\varepsilon^2)\)、\(I=(\varepsilon)\)。证明包含 \(I\hookrightarrow R\) 与 \(R/I\) 张量后变成来源非零的零映射，并求出来源和目标。
\end{exercise}
\begin{solution}
由\cref{b2:prop:tensor-quotient}，\(I\otimes_R(R/I)\cong I/I^2=I\)，它是一维 \(k\) 空间；目标 \(R\otimes_R(R/I)\cong R/I\cong k\)。包含诱导的映射把 \(\varepsilon\otimes\bar1\) 送到 \(\varepsilon\otimes\bar1=1\otimes\bar\varepsilon=0\)，所以为零。这里原包含单射，失效来自张量后的新关系 \(\bar\varepsilon=0\)，并不是来源原来就为零。
\end{solution}

\begin{exercise}\label{b2:ex:06-split-tensor}
设 \(0\to A\xrightarrow{i}B\xrightarrow{q}C\to0\) 是分裂的左 \(R\) 模短正合列。给定截面 \(s:C\to B\)，对任意右 \(R\) 模 \(M\)，证明张量后的序列仍分裂，并写出所用的截面。说明为什么这不与整数乘 \(n\) 的反例矛盾。
\end{exercise}
\begin{solution}
由\cref{b2:thm:split-exact-sequence}，存在 \(r:B\to A\) 满足 \(ri=\operatorname{id}_A\)，且 \(qs=\operatorname{id}_C\)。函子性给出张量后的缩回 \(\operatorname{id}_M\otimes r\) 与截面 \(\operatorname{id}_M\otimes s\)。缩回保证第一箭头单射，右正合性保证其余正合，截面保证分裂。整数序列 \(0\to\mathbb Z\xrightarrow{\times n}\mathbb Z\to\mathbb Z/n\mathbb Z\to0\) 不分裂：任何 \(\mathbb Z/n\mathbb Z\to\mathbb Z\) 的同态都是零，不可能成为商映射的截面。
\end{solution}

\begin{exercise}\label{b2:ex:06-localization-calculation}
计算 \(\mathbb Z[1/6]\otimes_{\mathbb Z}(\mathbb Z/40\mathbb Z)\) 和 \(\mathbb Q\otimes_{\mathbb Z}(\mathbb Z^2\oplus\mathbb Z/40\mathbb Z)\)。在第一个例子中写出分式到所得循环群的显式映射。
\end{exercise}
\begin{solution}
第一个张量积是 \(\Sigma^{-1}(\mathbb Z/40\mathbb Z)\)，其中 \(\Sigma=\{6^j:j\ge0\}\)。定义 \(\bar a/6^j\mapsto\bar a\in\mathbb Z/5\mathbb Z\)，因为 \(6\equiv1\pmod5\)。若分式相等，某个 \(6^u\) 零化交叉差，模 \(5\) 后即可消去它，说明映射良定义；它满射。若像为零，\(a\) 被 \(5\) 整除，因 \(8\mid6^3\)，有 \(40\mid6^3a\)，由\eqref{b2:eq:localization-zero}，该分式为零。因此第一项为 \(\mathbb Z/5\mathbb Z\)。第二项由直和相容性与单位同构等于 \(\mathbb Q^2\)：挠分量中每个元素被非零整数 \(40\) 零化，所以它的有理化为零。
\end{solution}

\begin{exercise}\label{b2:ex:06-localization-injective}
设 \(R\) 为交换环，\(\Sigma\) 为乘法闭集，\(K_1,K_2\) 为 \(R\) 模 \(N\) 的子模。将它们的局部化按\cref{b2:prop:localization-exact}视为 \(\Sigma^{-1}N\) 的子模，证明
\(\Sigma^{-1}(K_1\cap K_2)=\Sigma^{-1}K_1\cap\Sigma^{-1}K_2\)。
再取 \(N=\mathbb Z\)、\(\Sigma=\{2^j:j\ge0\}\)、\(K_i=2^i\mathbb Z\)（\(i\ge1\)），说明局部化不一定保持无限交。解释有限交的证明在哪一步不能直接用于无限交。
\end{exercise}
\begin{solution}
先说明 \(n/s\in\Sigma^{-1}K\) 当且仅当存在 \(u\in\Sigma\) 使 \(un\in K\)。若 \(n/s=k/t\)、\(k\in K\)，则某个 \(v\in\Sigma\) 满足 \(v(tn-sk)=0\)，故 \((vt)n=vsk\in K\)。反过来，若 \(un\in K\)，则 \(n/s=(un)/(us)\) 属于 \(\Sigma^{-1}K\)。

交中的元素 \(n/s\) 分别给出 \(u_1n\in K_1\)、\(u_2n\in K_2\)。于是 \(u_1u_2n\in K_1\cap K_2\)，由上述判据属于左端。另一包含直接成立。

例中 \(\bigcap_{i\ge1}2^i\mathbb Z=0\)，其局部化为零；可是 \(2^i\) 在 \(\mathbb Z[1/2]\) 中可逆，每个 \(\Sigma^{-1}K_i=\mathbb Z[1/2]\)，故这些局部化的交为 \(\mathbb Z[1/2]\)。有限交只需将有限多个 \(u_i\) 相乘，无限交则未必存在一个共同的 \(u\)。
\end{solution}

\begin{exercise}\label{b2:ex:06-field-extension-matrix}
令 \(V=\mathbb Q^2\)，\(T=\begin{pmatrix}0&2\\1&0\end{pmatrix}\)，\(K=\mathbb Q(\sqrt2)\)。求 \(K\otimes_{\mathbb Q}V\) 的一组 \(K\) 基及扩张后算子的特征向量。解释扩张前后维数以及是否可对角化的变化。
\end{exercise}
\begin{solution}
基为 \(1\otimes e_1,1\otimes e_2\)，所以 \(K\) 维数仍为 \(2\)。在此基下矩阵仍为 \(T\)。向量 \((\sqrt2,1)\)、\((-\sqrt2,1)\) 的特征值分别是 \(\sqrt2,-\sqrt2\)，它们的行列式为 \(2\sqrt2\ne0\)，故构成特征基。原算子的特征多项式 \(t^2-2\) 在 \(\mathbb Q\) 中无根，所以在 \(\mathbb Q\) 上不可对角化。若把扩张后的空间再限制成 \(\mathbb Q\) 空间，其维数为 \(4\)；“维数仍为 \(2\)”所指的是新基域 \(K\) 上的维数。
\end{solution}

\begin{exercise}\label{b2:ex:06-iterated-extension}
对环同态 \(R\xrightarrow{\alpha}S\xrightarrow{\varphi}T\) 与左 \(R\) 模 \(M\)，证明
\[
T\otimes_S(S\otimes_RM)\cong T\otimes_RM,
\qquad t\otimes(s\otimes m)\longmapsto t\cdot\varphi(s)\otimes m,
\]
给出逆映射并核对两个环上的平衡关系。
\end{exercise}
\begin{solution}
关于 \(S\) 的平衡性来自 \(t\cdot\varphi(s')\varphi(s)=t\cdot\varphi(s's)\)，关于 \(R\) 的平衡性来自
\[
t\cdot\varphi(s)\varphi\bigl(\alpha(r)\bigr)\otimes m
=t\cdot\varphi(s)\otimes rm.
\]
因此公式诱导左 \(T\) 线性映射。逆为 \(t\otimes m\mapsto t\otimes(1\otimes m)\)；若把 \(r\) 从 \(t\) 移向 \(m\)，先用外层 \(S\) 平衡关系，再用内层 \(R\) 平衡关系，即得同样的像。两复合由 \(t\otimes(s\otimes m)=t\cdot\varphi(s)\otimes(1\otimes m)\) 与单位作用等式成为恒等。这也可看作\cref{b2:thm:tensor-associativity-unit}的结合与单位同构的复合。
\end{solution}
